\documentclass[11pt,oneside]{report}
\usepackage{nm1notes}

\begin{document}

\begin{titlepage}
  \raggedright
  \vspace*{2in}
  {\Huge\bfseries CSCI-GA 2420 / MATH-GA 2010\par}
  \vspace{8pt}
  {\LARGE Numerical Methods I\par}
  \vspace{8pt}
  {\LARGE Course Notes\par}
  \vspace{8pt}
  {\Large Fall 2026\par}
  \vspace{24pt}
  {\large Instructor: Prof.\ Florian Schaefer\par}
  \vspace{4pt}
  {\large Notes by You Li\par}
\end{titlepage}

\thispagestyle{fancy}
\markright{}
{\Large\bfseries Acknowledgements\par}
\vspace{6pt}
These notes follow the Fall 2026 lectures of CSCI-GA 2420 / MATH-GA 2010 (Numerical Methods I) at the Courant Institute, taught by Prof.\ Florian Schaefer. Each chapter follows one lecture. Sections marked (SUPPLEMENT) go beyond the blackboard and draw on Trefethen and Bau~\cite{TB}, Demmel~\cite{Demmel}, and Golub and Van Loan~\cite{GVL}.

\paragraph{Notation.} Throughout, $\K$ denotes either $\R$ or $\C$. Vectors carry an arrow, $\vec{x} \in \K^n$; for $A \in \K^{m \times n}$, $A^\ast$ is the conjugate transpose (equal to $A^\T$ when $\K = \R$), $\range(A)$ is the range (column space), and $\nullsp(A)$ is the null space. The symbol $\defeq$ means ``is defined as''. Unless stated otherwise, $\norm{\cdot}$ is the Euclidean norm on vectors and the induced $2$-norm on matrices. Absolute values and inequalities between matrices, such as $\abs{E} \le \abs{L}\abs{U}$, are taken entrywise.

\tableofcontents

\chapter{Floating-Point Arithmetic, Conditioning, and Stability}
\label{ch:fp}

Supplementary reading:
\begin{itemize}
  \item \cite{TB} Lectures 12, 13, and 14.
  \item \cite{Demmel} Sections 1.3--1.5.
  \item \cite{GVL} Section 2.4.
\end{itemize}

\section{Why Exact Linear Algebra Is Not Enough}

In a linear algebra course, $\R$ is a field: addition is associative, multiplication distributes over addition, and $x$ and $(x + a) - a$ are the same number. Every derivation in that course relies on these facts.

A computer, however, has no $\R$. Its native objects are bits, $\B \doteq \{0, 1\}$, and any number system it offers must be built from finitely many of them. Approximating an uncountable set by a finite one causes two kinds of trouble:
\begin{enumerate}
  \item \emph{Representation error}: the number $x$ itself cannot be stored, so we store a nearby $\hat{x}$ instead.
  \item \emph{Propagated error}: even if the inputs are exact, the result of every operation has to be pushed back into the finite set.
\end{enumerate}
The second kind has a consequence that is worse than it first appears: \emph{floating-point addition is not associative}. Two algorithms that are mathematically equivalent can therefore return completely different answers on a computer. The question ``is this algorithm correct?'' splits off from the question ``is this algorithm numerically good?'', and this lecture builds the language for the second one. The \emph{condition number} measures how sensitive the problem itself is, and \emph{stability} measures whether an algorithm amplifies errors beyond that.

The topics of the semester are: numerical representation of vectors and matrices; solving linear systems by Gaussian elimination; solving least squares problems by orthogonalization; eigendecompositions; low-rank approximation; iterative methods; and structured matrices.

\section{Vectors and Matrices}

We write $\K$ for either $\R$ or $\C$, so that every result only needs to be stated once. A \emph{vector} $\vec{v} \in \K^m$ is an $m$-tuple of elements of $\K$, representing a position or direction in an $m$-dimensional coordinate system. A \emph{matrix} $A \in \K^{m \times n}$ is an $m \times n$ array of elements of $\K$, which may be thought of as $n$ positions in an $m$-dimensional coordinate system. We identify $\K^{m \times 1} \cong \K^m$, so a column vector is a matrix with one column:
\begin{equation}
  A = \begin{bmatrix} A_{11} & A_{12} \\ A_{21} & A_{22} \\ A_{31} & A_{32} \end{bmatrix} = \begin{bmatrix} \vec{a}_1 & \vec{a}_2 \end{bmatrix}.
\end{equation}
Scalar multiplication and addition act entrywise. For $A \in \K^{m \times n}$, $\vec{v} \in \K^n$, and $B \in \K^{n \times \ell}$,
\begin{equation}
  (A\vec{v})_i = \sum_{k} A_{ik} v_k, \qquad (AB)_{ij} = \sum_k A_{ik} B_{kj},
\end{equation}
so $A$ is a linear map $\K^n \to \K^m$ and $\K^{n \times \ell} \to \K^{m \times \ell}$: it carries a region, together with the vectors $\vec{x}$ in it, to the image vectors $A\vec{x}$.

The entrywise formula is rarely the most useful way to think. Numerical linear algebra relies much more on the \emph{column view}:
\begin{equation}
  A\vec{v} = \sum_{j=1}^{n} v_j \vec{a}_j,
\end{equation}
that is, $A\vec{v}$ is the linear combination of the columns of $A$ with coefficients given by the entries of $\vec{v}$. This view appears throughout the course. The column space $\range(A)$ decides whether $A\vec{x} = \vec{b}$ is solvable, the QR factorization orthogonalizes the columns, and low-rank approximation asks how few directions suffice to approximately span them. Likewise, the $j$-th column of $AB$ is $A$ applied to the $j$-th column of $B$, so a matrix--matrix product is $\ell$ matrix--vector products done at once.

\section{Representing Real Numbers}

We build up to floating-point numbers in four steps, each fixing a shortcoming of the previous one.

\subsection{Unsigned and Signed Integers}

A string $d \in \B^b$ represents the integers $0$ through $2^b - 1$ via
\begin{equation}
  \operatorname{UInt}(d) \doteq \sum_{i=1}^{b} d_i 2^{i-1}.
\end{equation}
To represent negative numbers, we use an offset and give the highest bit a negative weight:
\begin{equation}
  \operatorname{Int}(d) \doteq \sum_{i=1}^{b-1} d_i 2^{i-1} - d_b 2^{b-1},
\end{equation}
which covers $-2^{b-1}$ through $2^{b-1} - 1$. This is \emph{two's complement}. Nonnegative numbers keep their unsigned encoding and the adder circuit is shared with unsigned integers. The price is an asymmetric range: the most negative number, such as \texttt{typemin(Int8)} $= -128$, has no positive counterpart, so $-x$ can overflow back to $x$. (A different kind of offset, the \emph{biased} encoding that stores $x + K$, is used for floating-point exponents; see \Cref{sec:ieee}.)

\subsection{Fixed-Point Numbers}

Fixing the position of the binary point and using $f = \lfloor b/c \rfloor$ of the bits as fractional digits gives
\begin{equation}
  \operatorname{Fixed}(d) \doteq \operatorname{Int}(d) \cdot 2^{-f}.
\end{equation}
Fixed-point numbers have \emph{constant absolute spacing} $2^{-f}$ over the whole range. To represent numbers of size $s$ and $s^{-1}$ at the same time, one needs about $\log_2(s)$ bits, so precision and dynamic range are tied together and the cost grows \emph{linearly} in the dynamic range. Physical simulations routinely span a dozen orders of magnitude, which quickly breaks fixed-point arithmetic.

\subsection{Floating-Point Numbers}

The idea of floating point is to store the order of magnitude separately as an exponent. Given $(p, e) \in \B^{b_p} \times \B^{b_e}$, define
\begin{equation}
  \operatorname{Float}(p, e) \doteq (-1)^{p_1} \left(1 + \sum_{i=2}^{b_p} p_i 2^{-(i-1)}\right) 2^{\operatorname{Int}(e)}.
\end{equation}
The three parts are the sign bit $p_1$, the \emph{significand} (or \emph{mantissa}) in $[1, 2)$, and the power of two. The leading $1$ of the significand is \emph{implicit} and takes no storage. The decimal analogue is $\pm 1.123 \cdot 10^{-4}$.

The key benefit is that covering the range from $s^{-1}$ to $s$ only needs about $\log_2(\log_2(s))$ exponent bits. Fixed point pays a linear price; floating point pays a doubly logarithmic one. This is why floating point dominates scientific computing. A complex number $z \in \C$ is stored as two floating-point numbers. For details, the lecture recommends \cite{Overton}.

\section{The Geometry of Floating-Point Numbers}

Plotting all floating-point numbers on the real line shows that they are \emph{not equally spaced}. Within each \emph{binade} $[2^E, 2^{E+1})$ they are equally spaced, but crossing a power of two doubles the spacing: if the spacing near $1$ is $\eps \doteq 2^{-(b_p - 1)}$, it is $2\eps$ near $2$ and $4\eps$ near $4$.

\begin{center}
\begin{tikzpicture}[x=1.6cm]
  \draw[-{Latex}] (0.8,0) -- (8.6,0);
  \foreach \x in {1,1.25,...,2} \draw (\x,0.12) -- (\x,-0.12);
  \foreach \x in {2,2.5,...,4} \draw (\x,0.12) -- (\x,-0.12);
  \foreach \x in {4,5,...,8} \draw (\x,0.12) -- (\x,-0.12);
  \node[below] at (1,-0.15) {$1$};
  \node[below] at (2,-0.15) {$2$};
  \node[below] at (4,-0.15) {$4$};
  \node[below] at (8,-0.15) {$8$};
  \draw[decorate,decoration={brace,amplitude=4pt}] (1,0.2) -- (1.25,0.2) node[midway,above=4pt] {$\eps$};
  \draw[decorate,decoration={brace,amplitude=4pt}] (2,0.2) -- (2.5,0.2) node[midway,above=4pt] {$2\eps$};
  \draw[decorate,decoration={brace,amplitude=4pt}] (4,0.2) -- (5,0.2) node[midway,above=4pt] {$4\eps$};
\end{tikzpicture}
\end{center}

So although the absolute spacing changes, the \emph{relative spacing is nearly constant}: within one binade it lies between $\eps/2$ and $\eps$. This determines the philosophy behind floating point: it aims for \emph{relative} accuracy, on the premise that every quantity carries an uncertainty proportional to its own size.

\begin{example}[Floating-Point Addition Is Not Associative]
  \label{ex:nonassoc}
  Let $\eps$ be the spacing at $1$ and suppose ties are rounded up. Evaluate $(1 + \eps) + 1 + 2$ in two ways.
  \begin{itemize}
    \item Left to right: $(1 + \eps) + 1 = 2 + \eps$ lies exactly halfway between $2$ and $2 + 2\eps$ and rounds up to $2 + 2\eps$. Adding $2$ gives $4 + 2\eps$, halfway between $4$ and $4 + 4\eps$, which rounds up to $4 + 4\eps$.
    \item Last two terms first: $1 + 2 = 3$ exactly, and $(1 + \eps) + 3 = 4 + \eps$ rounds down to $4$.
  \end{itemize}
  The same expression yields $4$ or $4 + 4\eps$. Multiplication is not associative either, and the distributive law fails as well; commutativity of addition and of multiplication does hold.
\end{example}

The practical consequences are that a parallel reduction depends on the number of threads, that compiler flags such as \texttt{-ffast-math} may reorder sums and change the result, and that bitwise reproducibility across machines takes deliberate effort.

\section{Cancellation}

\begin{example}[Cancellation Exposes Existing Error]
  \label{ex:cancel}
  Suppose $10^6 + 1$ already carries a relative error $\eps$, and we subtract $10^6$:
  \begin{equation}
    (10^6 + 1)(1 + \eps) - 10^6 = 1 + 10^6 \eps.
  \end{equation}
  The exact answer is $1$ and the error is $10^6\eps$. The absolute error is unchanged, but the result is six orders of magnitude smaller than the operands, so the relative error has grown by a factor of $10^6$.
\end{example}

The subtraction itself is innocent: it introduces only $O(\eps)$ new error, and when the operands are close it introduces none at all (see Sterbenz's lemma in \Cref{sec:fp-further}). Cancellation merely \emph{exposes} error that was already present. A more accurate version of the slogan ``avoid cancellation'' is therefore: do not compute a small result as the difference of two large numbers.

\section{Conditioning}
\label{sec:conditioning}

How do errors \emph{propagate}? Consider computing $y = f(x)$ for $f : \R \to \R$. Since the input is rounded before anything else happens, the best we can hope for is $\hat{y} = f(\hat{x})$ with $\abs{\hat{x} - x}/\abs{x} < \eps$. Comparing relative errors,
\begin{equation}
  \frac{\abs{\hat{y} - y}/\abs{y}}{\abs{\hat{x} - x}/\abs{x}} = \frac{\abs{f(\hat{x}) - f(x)}}{\abs{\hat{x} - x}} \cdot \frac{\abs{x}}{\abs{f(x)}} \xrightarrow[\eps \to 0]{} \frac{\abs{x}\abs{f'(x)}}{\abs{f(x)}}.
\end{equation}

\begin{definition}[Condition Number]
  \label{def:cond}
  The (relative) \emph{condition number} of $f$ at $x$ is
  \begin{equation}
    \kappa_f(x) \doteq \frac{\abs{x}\abs{f'(x)}}{\abs{f(x)}}.
  \end{equation}
  The problem is \emph{well-conditioned} at $x$ if $\kappa_f(x) \approx 1$ and \emph{ill-conditioned} if $\kappa_f(x) \gg 1$.
\end{definition}

The condition number depends only on the problem and the input, not on the algorithm. This is the most important point of the lecture: no algorithm can rescue an ill-conditioned problem, because the difficulty is intrinsic.

\begin{example}[Conditioning of Subtraction]
  Let $f(x) = x - 10^6$ at $x = 10^6 + 1$. Then
  \begin{equation}
    \kappa_f(10^6 + 1) = \frac{\abs{10^6 + 1} \cdot 1}{1} \approx 10^6 \gg 1.
  \end{equation}
  Cancellation is thus a conditioning phenomenon, not a defect of floating point. Even an algorithm with infinite precision turns a relative uncertainty of $10^{-16}$ in $x$ into a relative uncertainty of $10^{-10}$ in the output.
\end{example}

\section{Stability}
\label{sec:stability}

Conditioning is about the problem; stability is about the algorithm. Let $\hat{f}$ be the floating-point algorithm we actually run: it acts on $\hat{x}$ and rounds at every internal step.

\begin{definition}[Stability]
  An algorithm $\hat{f}$ for $f$ is \emph{stable} if
  \begin{equation}
    \frac{\abs{\hat{f}(\hat{x}) - y}/\abs{y}}{\abs{\hat{x} - x}/\abs{x}} \lesssim \kappa_f(x),
  \end{equation}
  that is, if it does not amplify input perturbations more than the problem itself does. Otherwise $\hat{f}$ is \emph{unstable}.
\end{definition}

\begin{example}[An Unstable Algorithm for the Identity]
  \label{ex:unstable-identity}
  Let $f(x) = x$ and $\hat{f}(x) = (x + 10^6) - 10^6$. Mathematically $\hat{f} = f$, and the identity has $\kappa_f \equiv 1$, the best possible conditioning. But $\hat{f}$ first adds $x$ to $10^6$, which pushes the low-order bits of $x$ out of the significand for good, and then subtracts $10^6$ again. For $x = 1$ the relative error is $O(10^6 \eps)$ instead of $O(\eps)$.

  The diagnosis: \emph{the algorithm splits a well-conditioned problem into ill-conditioned intermediate steps}. Writing $\hat{f} = g_2 \circ g_1$ with $g_1(x) = x + 10^6$ and $g_2(z) = z - 10^6$, the step $g_1$ is well-conditioned but $g_2$ has condition number about $10^6$ at $z \approx 10^6$. A small overall condition number does not mean every step is well-conditioned, and the algorithm's error accumulates step by step. This view in terms of intermediate condition numbers will return when we compare Gram--Schmidt with the normal equations.
\end{example}

``Unstable'' and ``ill-conditioned'' are different diagnoses with different remedies: the former calls for a different algorithm, the latter for a different formulation of the problem (reparametrization, regularization, or higher precision).

\section{(SUPPLEMENT) IEEE 754 Formats}
\label{sec:ieee}

The formats in common use are listed below; $b_p$ counts the implicit bit, so a significand field with $b_p - 1$ stored bits provides $b_p$ bits of precision.
\begin{center}
\begin{tabular}{lcccccc}
  \toprule
  Format & Total bits & Exponent bits & Stored significand bits & $b_p$ & Bias & Decimal digits \\
  \midrule
  Float64 & 64 & 11 & 52 & 53 & 1023 & $\approx 15.9$ \\
  Float32 & 32 & 8 & 23 & 24 & 127 & $\approx 7.2$ \\
  Float16 & 16 & 5 & 10 & 11 & 15 & $\approx 3.3$ \\
  BFloat16 & 16 & 8 & 7 & 8 & 127 & $\approx 2.4$ \\
  \bottomrule
\end{tabular}
\end{center}
The exponent field uses a biased encoding, and its two extreme values are reserved:
\begin{itemize}
  \item An all-zeros exponent encodes $\pm 0$ (if the significand is also zero) or a \emph{subnormal} number, which drops the implicit leading $1$ and uses $0.f \times 2^{E_{\min}}$. Subnormals provide \emph{gradual underflow}, which guarantees $x \ne y \implies x - y \ne 0$.
  \item An all-ones exponent encodes $\pm\infty$ (zero significand) or \texttt{NaN} (nonzero significand).
\end{itemize}
The normal range of Float64 is therefore roughly $2.2 \times 10^{-308}$ to $1.8 \times 10^{308}$. IEEE defines five exceptions (invalid, division by zero, overflow, underflow, inexact), and specifies that $\infty$ and \texttt{NaN} propagate through arithmetic instead of halting it. Demmel~\cite[\S1.5]{Demmel} stresses the value of this: an algorithm can run to completion and check its result at the end, instead of branching at every step. Note that \texttt{NaN != NaN} is always true, so one must test with \texttt{isnan}.

\section{(SUPPLEMENT) Machine Epsilon and the Model of Arithmetic}

\subsection{Two Conventions for Machine Epsilon}

Two different quantities are both called ``machine epsilon'' in the literature, and they differ by a factor of $2$.

\begin{definition}[Machine Epsilon and Unit Roundoff]
  \label{def:eps}
  Let $b_p$ be the precision of a floating-point format.
  \begin{itemize}
    \item The \emph{machine epsilon} $\eps = \epsm \doteq 2^{-(b_p - 1)}$ is the distance from $1$ to the next larger floating-point number. For Float64, $\eps = 2^{-52} \approx 2.22 \times 10^{-16}$. This is the convention of the lecture, of Julia's \texttt{eps(Float64)}, and of C's \texttt{DBL\_EPSILON}.
    \item The \emph{unit roundoff} $u \doteq 2^{-b_p} = \eps/2$ is the largest relative rounding error, i.e., half a spacing. For Float64, $u = 2^{-53} \approx 1.11 \times 10^{-16}$. Trefethen and Bau~\cite{TB} call this quantity $\epsilon_{\text{machine}}$, and Golub and Van Loan~\cite{GVL} call it the unit roundoff.
  \end{itemize}
\end{definition}

Any statement of the form $O(\eps)$ is unaffected by the choice, but constants copied from different textbooks must be read with care.

The IEEE default rounding mode is \emph{round to nearest, ties to even}, under which both orders in \Cref{ex:nonassoc} give $4$. Non-associativity does not depend on the rounding mode, however. With ties to even,
\begin{equation}
  (1 + u) + u = 1 \qquad \text{but} \qquad 1 + (u + u) = 1 + \eps,
\end{equation}
since on the left each intermediate result is a tie that rounds back to $1$, while on the right $u + u = \eps$ and $1 + \eps$ are both exact.

\subsection{The Fundamental Axiom}

Error analysis needs rounding to be captured by a rule we can reason with. Let $\fl(x)$ denote the floating-point number nearest to $x$. Because the relative spacing is bounded, every $x$ in the normal range satisfies
\begin{equation}
  \fl(x) = x(1 + \delta), \qquad \abs{\delta} \le u. \label{eq:fl}
\end{equation}
IEEE requires $+$, $-$, $\times$, $\div$, and $\sqrt{\cdot}$ to be \emph{correctly rounded}: the exact result is computed in infinite precision and then rounded once. This gives the following model, stated in the same form in \cite[Lecture 13]{TB}, \cite[\S2.4]{GVL}, and \cite[\S1.5]{Demmel}.

\begin{theorem}[Fundamental Axiom of Floating-Point Arithmetic]
  \label{thm:axiom}
  For all floating-point numbers $x, y$ and $\circledast \in \{+, -, \times, \div\}$,
  \begin{equation}
    \fl(x \circledast y) = (x \circledast y)(1 + \delta), \qquad \abs{\delta} \le u.
  \end{equation}
\end{theorem}

The axiom has an immediate \emph{backward} reading. For addition,
\begin{equation}
  \fl(x + y) = (x + y)(1 + \delta) = x(1 + \delta) + y(1 + \delta),
\end{equation}
so the computed sum is \emph{exactly} the sum of two slightly perturbed inputs. Pushing errors back onto the inputs in this way is the main thread of error analysis in this course; see \Cref{sec:backward}.

\subsection{Accumulation of Errors}

A single operation incurs an error of size $O(u)$. For $n$ operations, Golub and Van Loan write
\begin{equation}
  \gamma_n \doteq \frac{nu}{1 - nu} \approx nu \qquad (nu \ll 1).
\end{equation}

\begin{proposition}[Error in Inner Products]
  \label{prop:dot}
  For $\vec{x}, \vec{y} \in \R^n$,
  \begin{equation}
    \abs{\fl(\vec{x}^\T \vec{y}) - \vec{x}^\T \vec{y}} \le \gamma_n \sum_{i=1}^{n} \abs{x_i}\abs{y_i} = \gamma_n \abs{\vec{x}}^\T \abs{\vec{y}}.
  \end{equation}
  Similarly, for matrices, $\abs{\fl(AB) - AB} \le \gamma_n \abs{A}\abs{B}$ entrywise.
\end{proposition}

The \emph{shape} of this bound matters much more than its constant. The right-hand side involves $\abs{\vec{x}}^\T\abs{\vec{y}}$, not $\abs{\vec{x}^\T\vec{y}}$. If all terms have the same sign, the two agree and the relative error is $O(nu)$, which is fine. If instead $\vec{x}^\T\vec{y}$ is a small number obtained by cancelling many terms of both signs, then $\abs{\vec{x}}^\T\abs{\vec{y}} \gg \abs{\vec{x}^\T\vec{y}}$ and the relative error can be arbitrarily bad. The linear factor $n$ is almost never the real issue (it is a very pessimistic worst case, and random errors grow more like $\sqrt{n}$). The real issue is always \emph{cancellation}.

\section{(SUPPLEMENT) Backward Stability}
\label{sec:backward}

The textbooks~\cite[Lecture 14]{TB} mostly use a stronger notion than the stability of \Cref{sec:stability}.

\begin{definition}[Backward Stability]
  \label{def:backward-stable}
  An algorithm $\hat{f}$ for $f$ is \emph{backward stable} if for every $x$ there exists $\tilde{x}$ with
  \begin{equation}
    \hat{f}(x) = f(\tilde{x}) \qquad \text{and} \qquad \frac{\abs{\tilde{x} - x}}{\abs{x}} = O(u).
  \end{equation}
\end{definition}

In words: \emph{the computed answer is the exact answer to a nearby problem}. The backward reading of \Cref{thm:axiom} is the simplest example. Backward stability is the gold standard of numerical linear algebra because it cleanly separates the problem from the algorithm, and it yields the classical rule of thumb
\begin{equation}
  \text{forward relative error} \lesssim \kappa_f(x) \cdot u,
\end{equation}
or, in a form that is easier to remember,
\begin{equation}
  \text{accuracy} \approx \text{condition number (problem)} \times \text{backward error (algorithm)}.
\end{equation}
A backward stable algorithm still returns inaccurate answers for ill-conditioned problems, and this is \emph{not the algorithm's fault}; it cannot be improved upon. Conversely, when a result is inaccurate, the first step should be to estimate the condition number, not to suspect the code.

\begin{center}
\begin{tabular}{l p{0.2\textwidth} p{0.34\textwidth} l}
  \toprule
  Concept & Belongs to & Question it answers & Measure \\
  \midrule
  Conditioning & problem $f$ and input $x$ & how much input perturbations are intrinsically amplified & $\kappa_f(x)$ \\
  Stability & algorithm $\hat{f}$ & whether the implementation amplifies errors further & backward error $O(u)$ \\
  Accuracy & final result & whether the computed number is right & $\lesssim \kappa \cdot u$ \\
  \bottomrule
\end{tabular}
\end{center}

Some common condition numbers:
\begin{center}
\begin{tabular}{lll}
  \toprule
  Problem & $\kappa$ & Remark \\
  \midrule
  $f(x) = \sqrt{x}$ & $1/2$ & always well-conditioned; even shrinks errors \\
  $f(x) = x^n$ & $n$ & grows linearly with the power \\
  $f(x) = e^x$ & $\abs{x}$ & ill-conditioned for large $\abs{x}$ \\
  $f(x_1, x_2) = x_1 - x_2$ & $\dfrac{\abs{x_1} + \abs{x_2}}{\abs{x_1 - x_2}}$ & cancellation (componentwise perturbations) \\
  $\vec{x} \mapsto A\vec{x}$ & $\dfrac{\norm{A}\norm{\vec{x}}}{\norm{A\vec{x}}} \le \kappa(A)$ & depends on the direction of $\vec{x}$ \\
  solve $A\vec{x} = \vec{b}$ & $\kappa(A) = \norm{A}\norm{A^{-1}}$ & \\
  polynomial roots & can be enormous & Wilkinson's polynomial \\
  \bottomrule
\end{tabular}
\end{center}
For $f : \R^n \to \R^m$ one replaces $f'$ by the Jacobian $J(\vec{x})$: the absolute condition number is $\hat{\kappa} = \norm{J(\vec{x})}$ and the relative condition number is $\kappa = \norm{J(\vec{x})}\norm{\vec{x}}/\norm{f(\vec{x})}$~\cite[Lecture 12]{TB}.

\section{(SUPPLEMENT) Further Topics}
\label{sec:fp-further}

\begin{strategy}[Rewriting to Avoid Cancellation]
  Common traps and their fixes:
  \begin{itemize}
    \item The root $\frac{-b + \sqrt{b^2 - 4ac}}{2a}$ of a quadratic cancels when $b^2 \gg 4ac$; compute the other root first and recover this one from Vieta's relation $x_1 x_2 = c/a$.
    \item The variance formula ``mean of squares minus square of mean'' cancels; use a two-pass algorithm or Welford's recurrence.
    \item $1 - \cos x$, $\sqrt{x + 1} - \sqrt{x}$, and $\log(1 + x)$ cancel as $x \to 0$; use a half-angle identity, rationalize, or call \texttt{log1p} and \texttt{expm1}.
  \end{itemize}
\end{strategy}

\begin{lemma}[Sterbenz]
  If $x, y$ are floating-point numbers with $y/2 \le x \le 2y$, then $x - y$ is a floating-point number, so $\fl(x - y) = x - y$ exactly.
\end{lemma}
This confirms again that cancellation does not lose new precision; it only exposes old errors.

\paragraph{Useful techniques.}
\begin{itemize}
  \item \emph{Kahan compensated summation} carries an extra variable that tracks the low-order bits lost in each rounding, reducing the error bound for a sum from $O(nu)$ to $O(u)$, independent of $n$. It costs four times as many operations, and \texttt{-ffast-math} may optimize it away.
  \item \emph{Fused multiply-add} (FMA) computes $a \cdot b + c$ with one rounding instead of two. It is a basic instruction on modern CPUs and GPUs, and it is why two mathematically equivalent expressions, or even \texttt{a*b - a*b}, may give different results.
  \item \texttt{-ffast-math} (or Julia's \texttt{@fastmath}) lets the compiler pretend that floating point is associative and distributive, trading correctness for speed. It should usually stay off in numerical code.
\end{itemize}

\paragraph{Low-precision formats.} Float16 and BFloat16 both use $16$ bits, but BFloat16 spends them on the exponent: it has $8$ exponent bits, the same range as Float32, and only $8$ bits of precision. This is the right trade-off for deep learning, where the first failure in training is usually overflow or underflow of gradients (a range problem) rather than lack of precision. Float16 has only $5$ exponent bits and a largest value of about $65504$, so training in Float16 requires loss scaling while BFloat16 does not. The accumulation bound $\gamma_n\abs{A}\abs{B}$ of \Cref{prop:dot} explains another convention, \emph{multiply in low precision, accumulate in high precision}: the factor $n$ comes from the length of the summation chain, and reduction dimensions of thousands are common, so tensor cores take BFloat16 or FP8 inputs but accumulate in Float32.

\paragraph{Experiments in Julia.}
\begin{lstlisting}
bitstring(1.0)                    # sign, exponent, and significand fields
eps(Float64), eps(1.0), eps(2.0)  # last two differ by 2: spacing doubles
nextfloat(1.0) - 1.0              # == eps(Float64) == 2^-52
0.1 + 0.2 == 0.3                  # false
(1.0 + eps()/2) + eps()/2 == 1.0 + (eps()/2 + eps()/2)   # false: not associative
1e16 + 1.0 - 1e16                 # compare with 1e6 + 1.0 - 1e6
Float16(65504), Float16(65520)    # the second overflows to Inf
nextfloat(Float64(0))             # smallest subnormal, about 5e-324
\end{lstlisting}

\chapter{Storage, Performance, and Gaussian Elimination}
\label{ch:performance}

Supplementary reading:
\begin{itemize}
  \item \cite{Demmel} Sections 2.3 and 2.6.
  \item \cite{TB} Lectures 17 and 20.
\end{itemize}

\section{Flops Are Not the Whole Story}

\Cref{ch:fp} was about accuracy: floating-point numbers have finite resolution, the condition number bounds the accuracy any algorithm can reach (a property of the problem), and stability measures whether an algorithm makes things worse (a property of the algorithm). Both concern whether we compute the \emph{right} answer.

Writing real numerical code runs into a second wall: how \emph{fast} we compute it. This wall has surprisingly little to do with the number of floating-point operations (flops). The lecture illustrates this with a Julia experiment timing $C \mathrel{+}= AB$ with $A \in \R^{\ell \times m}$, $B \in \R^{m \times n}$, for $\ell = 2 \cdot 10^4$, $m = 3 \cdot 10^4$ and varying $n$:
\begin{center}
\begin{tabular}{cccc}
  \toprule
  Size & Flops & Measured time & Throughput \\
  \midrule
  $n = 1$ & $6 \cdot 10^8$ & $6 \cdot 10^{-2}$ s & $\approx 10^{10}$ flop/s \\
  $n = 10^3$ & $6 \cdot 10^{11}$ & $4$ s & $\approx 1.5 \cdot 10^{11}$ flop/s \\
  \bottomrule
\end{tabular}
\end{center}
The flop count grows by a factor of $1000$, but the time grows only by a factor of about $70$: the cost per flop drops by more than an order of magnitude. When instead $\ell = m = n$ grows from $10^3$ ($7 \cdot 10^{-3}$ s) to $10^4$ ($6$ s), the throughput stays essentially constant, because the code already runs near the machine's peak.

The explanation is that an algorithm performs four kinds of operations whose costs are not remotely comparable:
\begin{equation}
  \operatorname{cost}(+) < \operatorname{cost}(\ast) \ll \operatorname{cost}([\ ]).
\end{equation}
In other words, \emph{fetching a number costs more than computing with it}. The case $n = 1$ is slow not because its multiply--adds are more expensive, but because its inner loop has length $1$: every number that is fetched is used once and discarded, so the code degenerates into pure data movement.

\section{Memory Is One-Dimensional}

Memory is a long string of bits, and a vector is a contiguous run of floating-point numbers $u_1, u_2, u_3, \dots$. There is no such thing as a two-dimensional matrix in memory; it must be flattened in some order. Julia, Fortran, BLAS, and LAPACK use \emph{column-major} order, storing entire columns contiguously:
\begin{equation}
  A_{11},\ A_{21},\ A_{31},\ A_{41},\ A_{12},\ A_{22},\ A_{32},\ A_{42},\ \dots
\end{equation}
so that $A_{ij}$ sits at linear index $i + (j - 1)m$. C and NumPy default to row-major order. Besides the data, one also stores metadata such as sizes, strides, and preallocated buffers.

Since the storage is just a contiguous block, the same memory can be \emph{reinterpreted} as an array of shape $m \times n \times k \times l$, $mn \times kl$, or $mnk \times l$ without any copying. Reshaping is free; transposing is not.

\section{The Memory Hierarchy}

Main memory is much slower than arithmetic, and hardware responds with caches: small, fast layers of memory close to the processor. A useful picture is a small warehouse next door (the cache) and a huge warehouse far away that must be reached by truck (RAM). Two different bottlenecks arise:
\begin{itemize}
  \item \emph{Latency}: how many cycles a single access waits. An L1 hit costs a few cycles; main memory costs two or three hundred.
  \item \emph{Bandwidth}: how many bytes can be moved per unit time.
\end{itemize}
Latency is fought with \emph{prefetching}: the hardware guesses what will be needed next and fetches it early. The guesses are good only when the access pattern is regular, i.e., contiguous or with a fixed stride. The \emph{regularity} of memory access is therefore often more important than its \emph{volume}. GPUs have no prefetcher in the traditional sense; instead, the threads of a warp (SIMT) execute the same instruction on consecutive elements, and the hardware coalesces their accesses into a single transaction. On a CPU contiguous access is faster; on a GPU it is close to mandatory.

\subsection{Loop Order Determines the Access Pattern}

\begin{example}[Two Loop Orders for a Matrix--Vector Product]
  The two loops below compute $\vec{u} = A\vec{v}$ with identical flop counts.
\begin{lstlisting}
# i outer: inner loop along a row of A       # j outer: inner loop down a column
for i = 1:m                                    for j = 1:n
    for j = 1:n                                    for i = 1:m
        u[i] += A[i,j] * v[j]                          u[i] += A[i,j] * v[j]
    end                                            end
end                                            end
\end{lstlisting}
  In column-major storage, the left version walks along a row of $A$ in the inner loop, jumping $m$ entries at each step and landing on a new cache line every time. The right version walks down a column of $A$ contiguously, and \texttt{u[i]} is contiguous as well; only \texttt{v[j]} changes in the outer loop. In the lecture's words, $j$-inside-$i$ has at least two random accesses, whereas $i$-inside-$j$ has only one.
\end{example}

The two versions are two different decompositions into vector operations. With $j$ inside, each row contributes an inner product (\texttt{dot}); with $i$ inside, the columns of $A$ are accumulated with weights $v_j$ (\texttt{axpy}). Column-major languages should use the \texttt{axpy} form and row-major languages the \texttt{dot} form. The lecture adds an important qualification: this matters most for \emph{small} matrices. For large matrices both versions are limited by bandwidth and the gap narrows; at small sizes such constant factors dominate.

\section{Tiling}

Matrix multiplication requires $\sim m^3$ flops but has only $m^2$ data. If $\approx m^2$ floating-point numbers fit into the cache, the $m^3$ flops can be performed after moving each number into the cache once. For larger matrices, we \emph{tile}: split $A$, $B$, $C$ into $b \times b$ blocks, $N = n/b$ blocks per side, and perform the multiplication block by block, choosing $b$ so that the blocks being worked on fit into the cache.

\begin{algorithm}[H]
  \caption{Tiled matrix multiplication $C \gets C + AB$.}
  \label{alg:tiled-matmul}
  \begin{algorithmic}[1]
    \For{$i = 1, \dots, N$}
      \For{$j = 1, \dots, N$}
        \State read block $C_{ij}$ into fast memory
        \For{$k = 1, \dots, N$}
          \State read blocks $A_{ik}$ and $B_{kj}$ into fast memory
          \State $C_{ij} \gets C_{ij} + A_{ik}B_{kj}$ \Comment{block product, entirely in fast memory}
        \EndFor
        \State write block $C_{ij}$ back to slow memory
      \EndFor
    \EndFor
  \end{algorithmic}
\end{algorithm}

\Cref{sec:intensity} quantifies how much this helps.

\section{Linear Systems}

Given $A \in \K^{m \times n}$ and $\vec{b} \in \range(A) \subseteq \K^m$, we seek $\vec{x} \in \K^n$ with $A\vec{x} = \vec{b}$. This \emph{can} be done by computing $A^{-1}$, but it need not and should not be: forming $A^{-1}$ costs $2n^3$ flops (versus $\tfrac{2}{3}n^3$ for the LU factorization below), still requires a matrix--vector product afterwards, and comes with a worse error bound.

For square $A \in \K^{m \times m}$, the following are equivalent:
\begin{equation}
  \nullsp(A) = \{\zerovec\} \iff \text{the columns of } A \text{ are linearly independent} \iff \range(A) = \K^m.
\end{equation}
The left condition says that $\vec{x} \mapsto A\vec{x}$ is injective (if $A\vec{x} = A\vec{y}$ then $A(\vec{x} - \vec{y}) = \zerovec$, so $\vec{x} = \vec{y}$), and the right condition says that it is surjective. We then call $A$ full rank, invertible, or nonsingular. The equivalence of injectivity and surjectivity holds only for square matrices; rectangular systems are treated in \Cref{ch:ls}.

\section{Triangular Systems}

We begin with the easy case. The small example
\begin{equation}
  \begin{aligned} x_1 + 5 &= 15 \\ x_1 + 5x_2 &= 20 \end{aligned} \qquad\Longrightarrow\qquad x_1 = 10,\ x_2 = 2
\end{equation}
already shows the whole idea: $x_1$ does not depend on $x_2$, so the triangular structure \emph{decouples} the unknowns and they can be solved for one at a time. For lower triangular $A$, \emph{forward substitution} computes
\begin{equation}
  x_1 = \frac{b_1}{A_{11}}, \qquad x_i = \frac{1}{A_{ii}}\left(b_i - \sum_{j=1}^{i-1} A_{ij}x_j\right), \quad i = 2, \dots, m.
\end{equation}
Upper triangular systems are solved in reverse order by \emph{back substitution}. The same formula admits two implementations, both costing $\sim m^2$ flops but with different access patterns.

\begin{algorithm}[H]
  \caption{Forward substitution, row-oriented (\texttt{dot} form).}
  \begin{algorithmic}[1]
    \For{$i = 1, \dots, m$}
      \State $x_i \gets b_i$
      \For{$j = 1, \dots, i - 1$}
        \State $x_i \gets x_i - A_{ij}x_j$
      \EndFor
      \State $x_i \gets x_i / A_{ii}$
    \EndFor
  \end{algorithmic}
\end{algorithm}

\begin{algorithm}[H]
  \caption{Forward substitution, column-oriented (\texttt{axpy} form), the lecture's efficient variant.}
  \begin{algorithmic}[1]
    \State $\vec{x} \gets \vec{b}$
    \For{$j = 1, \dots, m$}
      \State $x_j \gets x_j / A_{jj}$
      \For{$i = j + 1, \dots, m$}
        \State $x_i \gets x_i - A_{ij}x_j$ \Comment{remove the contribution of $x_j$ at once}
      \EndFor
    \EndFor
  \end{algorithmic}
\end{algorithm}

For back substitution, the loops run over $i = m, \dots, 1$ and $j = i + 1, \dots, m$ (row form), or over $j = m, \dots, 1$ and $i = 1, \dots, j - 1$ (column form). In column-major storage the row form walks along rows (strided) and the column form walks down columns (contiguous); this is the matrix--vector discussion again.

\section{Gaussian Elimination}

A general matrix is not triangular. Gaussian elimination \emph{transforms $A$ into upper triangular form by a sequence of simple operations}, reducing the general problem to a triangular one. A $2 \times 2$ example:
\begin{equation}
  \begin{aligned} x_1 + 2x_2 &= 10 \\ 2x_1 + x_2 &= 8 \end{aligned}
  \quad \xrightarrow{\ \text{row 2} \,-\, 2 \times \text{row 1}\ } \quad
  \begin{aligned} x_1 + 2x_2 &= 10 \\ -3x_2 &= -12 \end{aligned}
\end{equation}
In matrix language, this elimination step is a \emph{left multiplication}:
\begin{equation}
  \begin{bmatrix} 1 & 0 \\ -2 & 1 \end{bmatrix}
  \begin{bmatrix} 1 & 2 \\ 2 & 1 \end{bmatrix}
  \begin{bmatrix} x_1 \\ x_2 \end{bmatrix}
  =
  \begin{bmatrix} 1 & 0 \\ -2 & 1 \end{bmatrix}
  \begin{bmatrix} 10 \\ 8 \end{bmatrix}.
\end{equation}

\begin{theorem}[Elementary Lower Triangular Matrices]
  \label{thm:elim}
  Let $A \in \K^{m \times m}$ and suppose the first $k - 1$ columns of $A$ are already zero below the diagonal. Define the multipliers $\ell_{ik} \doteq A_{ik}/A_{kk}$ for $i = k + 1, \dots, m$, the vector
  \begin{equation}
    \vec{\ell}_k \doteq \begin{bmatrix} 0 & \cdots & 0 & \ell_{k+1,k} & \cdots & \ell_{mk} \end{bmatrix}^\T,
  \end{equation}
  whose first $k$ entries vanish, and the \emph{elementary lower triangular matrix}
  \begin{equation}
    L_k \doteq I - \vec{\ell}_k \vec{e}_k^{\,\ast} =
    \begin{bmatrix}
      1 & & & & \\
      & \ddots & & & \\
      & & 1 & & \\
      & & -\ell_{k+1,k} & \ddots & \\
      & & \vdots & & \\
      & & -\ell_{mk} & & 1
    \end{bmatrix}.
  \end{equation}
  Then $L_k A$ is zero below the diagonal in its first $k$ columns. Its first $k$ rows equal those of $A$, and only the trailing $(m - k) \times (m - k)$ block is updated.
\end{theorem}

\begin{proof}
  Left multiplication by $L_k$ performs, for each $i > k$, the row operation ``row $i$ minus $\ell_{ik}$ times row $k$''. The $(i, k)$ entry becomes $A_{ik} - \ell_{ik}A_{kk} = 0$. Row $k$ is zero in its first $k - 1$ columns, so the columns that have already been eliminated are not disturbed.
\end{proof}

Applying \Cref{thm:elim} for $k = 1, \dots, m - 1$ gives
\begin{equation}
  \underbrace{L_{m-1} \cdots L_1 A}_{\text{upper triangular}} = L_{m-1} \cdots L_1 \vec{b}.
\end{equation}
With $U \doteq L_{m-1} \cdots L_1 A$ and $L \doteq (L_{m-1} \cdots L_1)^{-1}$, we obtain $A = LU$ with $U$ upper triangular. The lecture stops here with two questions: \emph{is $L$ lower triangular, and can we compute it?} Both are answered in \Cref{ch:lu}.

\section{(SUPPLEMENT) Leading Dimension}

In column-major order, the entry $A_{ij}$ lives at
\begin{equation}
  \operatorname{addr}(A_{ij}) = \text{base} + (j - 1) \cdot \mathrm{lda} + (i - 1),
\end{equation}
where the \emph{leading dimension} $\mathrm{lda} \ge m$ is the distance between consecutive columns. Allowing $\mathrm{lda} \ne m$ lets us describe a \emph{submatrix} of a larger matrix by the quadruple (start address, $m$, $n$, $\mathrm{lda}$) without copying any data. This is what Julia's \texttt{view} and \texttt{@views} do, whereas \texttt{A[1:100, 1:100]} makes a copy.

\section{(SUPPLEMENT) Computational Intensity}
\label{sec:intensity}

The discussion of this lecture can be quantified~\cite[\S2.6]{Demmel}. Consider two levels of memory, a fast memory of size $M$ and a slow memory. Let one flop take time $t_f$ and one slow-memory access take time $t_m \gg t_f$.

\begin{definition}[Computational Intensity]
  If an algorithm performs $f$ flops and $m$ slow-memory accesses, its \emph{computational intensity} is $q \doteq f/m$.
\end{definition}

The total time is roughly
\begin{equation}
  T = f t_f + m t_m = f t_f \left(1 + \frac{t_m}{t_f} \cdot \frac{1}{q}\right).
\end{equation}
Running near the peak rate $f t_f$ requires $q \gg t_m / t_f$, and on modern machines $t_m / t_f$ is in the tens to hundreds. The three levels of the BLAS are distinguished precisely by $q$ (here $n$ is the vector length or matrix order):
\begin{center}
\begin{tabular}{lllll}
  \toprule
  Level & Typical operation & Flops $f$ & Accesses $m$ & $q$ \\
  \midrule
  Level 1 & \texttt{axpy} $\vec{y} \gets \alpha\vec{x} + \vec{y}$, \texttt{dot} & $2n$ & $3n$ & $2/3$ \\
  Level 2 & \texttt{gemv} $\vec{y} \gets \alpha A\vec{x} + \beta\vec{y}$, triangular solve & $2n^2$ & $n^2$ & $2$ \\
  Level 3 & \texttt{gemm} $C \gets \alpha AB + \beta C$ & $2n^3$ & $4n^2$ & $n/2$ \\
  \bottomrule
\end{tabular}
\end{center}
Only Level 3 has an intensity that grows with the problem size, which is the ``$m^3$ flops on $m^2$ data'' remark of the lecture. Triangular solves, in either loop order, are Level 2 with $q \approx 2$ and are fundamentally limited by bandwidth.

For naive matrix multiplication $C \mathrel{+}= AB$ with three $n \times n$ matrices, the innermost loop scans a whole column of $B$ each time. Then $A$ is read once ($n^2$), $C$ is read and written once ($2n^2$), and $B$ is read $n$ times ($n^3$), so
\begin{equation}
  q = \frac{2n^3}{n^3 + 3n^2} \approx 2,
\end{equation}
no better than Level 2. For the tiled version (\Cref{alg:tiled-matmul}) the slow-memory traffic is $m = 2n^2 + 2Nn^2$, so
\begin{equation}
  q \approx \frac{2n^3}{2Nn^2} = \frac{n}{N} = b.
\end{equation}
The intensity equals the block size. Three blocks must reside in fast memory simultaneously, so $3b^2 \le M$, i.e., $b \le \sqrt{M/3}$, and hence $q \approx \sqrt{M/3}$. A lower bound of Hong and Kung shows that any matrix multiplication that only reorders the sums has $q = O(\sqrt{M})$, so tiling is asymptotically optimal. Real implementations tile at \emph{several} levels (registers, L1, L2, and across nodes), each with its own $b$.

\section{(SUPPLEMENT) Blocked LU}
\label{sec:blocked-lu}

The in-place LU algorithm of \Cref{ch:lu} spends its time in rank-one updates, a Level 2 operation with $q \approx 2$, so naive Gaussian elimination is bandwidth-bound. The tiling idea fixes this~\cite[\S2.6.3]{Demmel}. Choose a block width $b$ and write
\begin{equation}
  A = \begin{bmatrix} A_{11} & A_{12} \\ A_{21} & A_{22} \end{bmatrix}, \qquad A_{11} \in \K^{b \times b}.
\end{equation}
\begin{enumerate}
  \item Factor the $m \times b$ panel $\begin{bmatrix} A_{11} \\ A_{21} \end{bmatrix}$ by ordinary elimination, obtaining $L_{11}$, $L_{21}$, $U_{11}$. This is still Level 2, but it accounts for only an $O(b/m)$ fraction of the work.
  \item Compute $U_{12} = L_{11}^{-1}A_{12}$, a triangular solve with many right-hand sides (Level 3).
  \item Compute $\tilde{A}_{22} = A_{22} - L_{21}U_{12}$, a matrix multiplication (Level 3) that contains most of the flops.
  \item Recurse on $\tilde{A}_{22}$.
\end{enumerate}
Only an $O(b/m)$ fraction of the $\tfrac{2}{3}m^3$ flops remains in bandwidth-bound Level 2 operations; the rest runs as matrix multiplication and enjoys $q \approx b$. This is where LAPACK's \texttt{getrf} gains its large speedup over the older LINPACK, and it is the first real instance of the rule ``rewrite everything as matrix multiplication''.

\chapter{Gaussian Elimination and the LU Factorization}
\label{ch:lu}

Supplementary reading:
\begin{itemize}
  \item \cite{TB} Lectures 20 and 21.
  \item \cite{Demmel} Section 2.3.
\end{itemize}

\section{Recap: From Triangular Systems to Elimination}

The lecture opens with a recap of \Cref{ch:performance}: matrices are stored column by column, memory access costs far more than arithmetic, and the loop order and tiling of an algorithm should match the storage. It then returns to linear systems. Triangular systems are solved by forward or back substitution in $O(m^2)$ flops, and Gaussian elimination reduces a general system to a triangular one: by \Cref{thm:elim}, applying elementary lower triangular matrices $L_k = I - \vec{\ell}_k\vec{e}_k^{\,\ast}$ for $k = 1, \dots, m - 1$ gives
\begin{equation}
  U \doteq L_{m-1} \cdots L_1 A \ \text{ upper triangular}, \qquad L \doteq (L_{m-1} \cdots L_1)^{-1}, \qquad A = LU.
\end{equation}
This lecture answers the two questions left open: \emph{is $L$ lower triangular, and can we compute it cheaply?} Answering them requires a little algebra of lower triangular matrices.

\section{The Algebra of Lower Triangular Matrices}

\begin{proposition}[Closure Properties]
  \label{prop:lower-closure}
  The lower triangular matrices in $\K^{m \times m}$ are closed under
  \begin{enumerate}
    \item addition,
    \item scalar multiplication,
    \item matrix multiplication, and
    \item inversion (of invertible ones).
  \end{enumerate}
\end{proposition}

\begin{proof}
  Items 1 and 2 are clear. For item 3, let $L$ and $\tilde{L}$ be lower triangular, so $(L\tilde{L})_{ij} = \sum_k L_{ik}\tilde{L}_{kj}$. A term $L_{ik}\tilde{L}_{kj}$ can be nonzero only if $i \ge k$ and $k \ge j$, hence $i \ge j$. So every term vanishes when $i < j$, and $L\tilde{L}$ is lower triangular. Item 4 is proved below for the matrices arising in LU, using \Cref{thm:elem-inverse}; the general case is in \Cref{sec:lu-supp}.
\end{proof}

\begin{theorem}[Inverses and Products of Elementary Lower Triangular Matrices]
  \label{thm:elem-inverse}
  Let $L_k = I - \vec{\ell}_k\vec{e}_k^{\,\ast}$ be as in \Cref{thm:elim}. Then:
  \begin{enumerate}
    \item $L_k^{-1} = I + \vec{\ell}_k\vec{e}_k^{\,\ast}$, i.e., inversion just \emph{flips the signs} of the multipliers;
    \item the product of the inverses, \emph{in this order}, places every multiplier in its own position:
    \begin{equation}
      L_1^{-1} L_2^{-1} \cdots L_{m-1}^{-1} =
      \begin{bmatrix}
        1 & & & \\
        \ell_{21} & 1 & & \\
        \vdots & \ddots & \ddots & \\
        \ell_{m1} & \cdots & \ell_{m,m-1} & 1
      \end{bmatrix}.
    \end{equation}
  \end{enumerate}
\end{theorem}

The order in item 2 matters: in the reverse order $L_{m-1}^{-1} \cdots L_1^{-1}$ the multipliers would contaminate each other. Both items rest on one observation: the first $k$ entries of $\vec{\ell}_k$ vanish, so $\vec{e}_p^{\,\ast}\vec{\ell}_k = 0$ for all $p \le k$.

\begin{proof}
  For item 1, since $\vec{e}_k^{\,\ast}\vec{\ell}_k = 0$,
  \begin{equation}
    (I + \vec{\ell}_k\vec{e}_k^{\,\ast})(I - \vec{\ell}_k\vec{e}_k^{\,\ast}) = I + \vec{\ell}_k\vec{e}_k^{\,\ast} - \vec{\ell}_k\vec{e}_k^{\,\ast} - \vec{\ell}_k(\vec{e}_k^{\,\ast}\vec{\ell}_k)\vec{e}_k^{\,\ast} = I.
  \end{equation}
  For item 2, compute each entry directly:
  \begin{equation}
    (L_1^{-1} \cdots L_{m-1}^{-1})_{ij} = \vec{e}_i^{\,\ast}(I + \vec{\ell}_1\vec{e}_1^{\,\ast}) \cdots (I + \vec{\ell}_{m-1}\vec{e}_{m-1}^{\,\ast})\vec{e}_j.
  \end{equation}
  Simplify from the right. For $p > j$, $\vec{e}_p^{\,\ast}\vec{e}_j = 0$, so $(I + \vec{\ell}_p\vec{e}_p^{\,\ast})\vec{e}_j = \vec{e}_j$ and those factors disappear; the $j$-th factor gives $(I + \vec{\ell}_j\vec{e}_j^{\,\ast})\vec{e}_j = \vec{e}_j + \vec{\ell}_j$:
  \begin{equation}
    = \vec{e}_i^{\,\ast}(I + \vec{\ell}_1\vec{e}_1^{\,\ast}) \cdots (I + \vec{\ell}_{j-1}\vec{e}_{j-1}^{\,\ast})(\vec{e}_j + \vec{\ell}_j).
  \end{equation}
  The part $\vec{e}_j$ passes through the remaining factors unchanged for the same reason, and so does $\vec{\ell}_j$, because $\vec{e}_p^{\,\ast}\vec{\ell}_j = 0$ for $p < j$. Hence the entry equals $\vec{e}_i^{\,\ast}\vec{\ell}_j + \vec{e}_i^{\,\ast}\vec{e}_j$, which is $1$ on the diagonal, $\ell_{ij}$ for $i > j$, and $0$ for $i < j$.
\end{proof}

\section{Answering the Two Questions}

\paragraph{$L$ is lower triangular.} We have
\begin{equation}
  L = (L_{m-1} \cdots L_1)^{-1} = L_1^{-1} \cdots L_{m-1}^{-1}.
\end{equation}
Each $L_k^{-1}$ is lower triangular by item 1 of \Cref{thm:elem-inverse}, and so is their product by item 3 of \Cref{prop:lower-closure}.

\paragraph{$L$ can be written down directly.} By item 2 of \Cref{thm:elem-inverse}, $L$ is unit lower triangular, with the multipliers $\ell_{ik}$ used during elimination sitting untouched in position $(i, k)$. No $L_k$ ever needs to be formed and no matrix product or inverse needs to be computed.

\section{The In-Place Algorithm}

Computing the LU factorization therefore amounts to \emph{transforming $A$ into $U$ while recording the multipliers}. The unit diagonal of $L$ need not be stored, and the strictly lower part of $L$ fits exactly into the zeros created in $U$. Hence $A$ can be overwritten by $L$ and $U$ with no extra storage.

\begin{algorithm}[H]
  \caption{LU factorization without pivoting, in place.}
  \label{alg:lu}
  \begin{algorithmic}[1]
    \For{$k = 1, \dots, m - 1$}
      \For{$i = k + 1, \dots, m$}
        \State $A_{ik} \gets A_{ik} / A_{kk}$ \Comment{multiplier $\ell_{ik}$, stored in place}
      \EndFor
      \For{$j = k + 1, \dots, m$}
        \For{$i = k + 1, \dots, m$}
          \State $A_{ij} \gets A_{ij} - A_{ik}A_{kj}$ \Comment{update the trailing block}
        \EndFor
      \EndFor
    \EndFor
  \end{algorithmic}
\end{algorithm}

On exit, the upper triangle of $A$ (including the diagonal) holds $U$ and the strictly lower triangle holds the multipliers of $L$. The innermost loop runs over $i$, i.e., down a column, matching the column-major storage as recommended in \Cref{ch:performance}.

\section{Zero Pivots}

The algorithm divides by $A_{kk}$, the $k$-th \emph{pivot}. If a pivot is zero, it cannot proceed.

\begin{example}[An Invertible Matrix with No LU Factorization]
  The matrix
  \begin{equation}
    A = \begin{bmatrix} 0 & 1 \\ 1 & 0 \end{bmatrix}
  \end{equation}
  is invertible, but $A_{11} = 0$ is a zero pivot, so the very first step divides by zero and $A = LU$ does not exist. \emph{Invertibility does not imply the existence of an LU factorization.}
\end{example}

\begin{theorem}[Zero Pivots and Leading Principal Submatrices]
  \label{thm:lu-exists}
  If the $k$-th pivot of the LU factorization of $A$ is zero, then the leading principal submatrix $A[1{:}k, 1{:}k]$ is singular. All pivots of the LU factorization are nonzero if and only if every leading principal submatrix of $A$ is nonsingular.
\end{theorem}

\begin{proof}
  From $A = LU$, look at the first $k$ rows and columns:
  \begin{equation}
    A[1{:}k, 1{:}k] = L[1{:}k, :]\,U[:, 1{:}k] = L[1{:}k, 1{:}k]\,U[1{:}k, 1{:}k].
  \end{equation}
  The second equality holds because $L$ is lower triangular, so $L[1{:}k, k{+}1{:}m] = 0$ and the rows of $U$ below row $k$ do not contribute. The matrix $L[1{:}k, 1{:}k]$ is unit lower triangular, hence nonsingular. The matrix $U[1{:}k, 1{:}k]$ is upper triangular with determinant equal to the product of the first $k$ pivots. Thus $A[1{:}k, 1{:}k]$ is nonsingular if and only if the first $k$ pivots are nonzero. When elimination reaches step $k$, the first $k$ rows and columns already satisfy the identity above, so the argument applies even halfway through the algorithm.
\end{proof}

\section{Small Pivots Are Also Dangerous}

Exactly zero pivots are unlikely in practice. Is it then enough that no pivot is exactly zero? Consider a small $\eps > 0$:
\begin{equation}
  A = \begin{bmatrix} \eps & 1 \\ 1 & 1 \end{bmatrix}, \qquad
  L = \begin{bmatrix} 1 & 0 \\ \eps^{-1} & 1 \end{bmatrix}, \qquad
  U = \begin{bmatrix} \eps & 1 \\ 0 & 1 - \eps^{-1} \end{bmatrix}.
\end{equation}
This is mathematically correct. In floating-point arithmetic, however, $\eps^{-1}$ is huge and $1 - \eps^{-1}$ rounds to $-\eps^{-1}$. The computed factors then give
\begin{equation}
  \tilde{L}\tilde{U} = \begin{bmatrix} \eps & 1 \\ 1 & \eps^{-1} - \eps^{-1} \end{bmatrix} = \begin{bmatrix} \eps & 1 \\ 1 & 0 \end{bmatrix} \approx \begin{bmatrix} 0 & 1 \\ 1 & 0 \end{bmatrix},
\end{equation}
which is nothing like $A$. The small pivot produces a huge multiplier, and rounding wipes out all information in $A_{22}$.

\section{Error Analysis of the LU Factorization}

\begin{theorem}[Backward Error of LU]
  \label{thm:lu-backward-preview}
  Let $\epsm$ be the machine epsilon. If the LU factorization runs to completion in floating-point arithmetic without encountering a zero pivot, the computed factors $L$, $U$ satisfy
  \begin{equation}
    A = LU + E, \qquad \abs{E} \le \epsm\, m\abs{L}\abs{U},
  \end{equation}
  where $\abs{\cdot}$ takes entrywise absolute values and the inequality holds entrywise.
\end{theorem}

The proof is given in \Cref{ch:pivoting}. The theorem says that the computed factorization is the exact factorization of a nearby matrix, with a perturbation controlled by $\abs{L}\abs{U}$. If $\abs{L}\abs{U}$ is comparable to $\abs{A}$, the error is at the level of machine precision. If $L$ or $U$ contains huge entries, the bound is useless. In the example above, $L$ contains $\eps^{-1}$ and $U$ contains $1 - \eps^{-1}$, so $\abs{L}\abs{U}$ has entries of size $\eps^{-1}$ and the error can be as large as $A$ itself.

\section{(SUPPLEMENT) Further Topics}
\label{sec:lu-supp}

\paragraph{Inverses of general lower triangular matrices.} The lecture proves item 4 of \Cref{prop:lower-closure} only for matrices of the form $(L_{m-1} \cdots L_1)^{-1}$. It holds for every invertible lower triangular $T$. Invertibility forces the diagonal entries to be nonzero; let $D = \diag(T)$, so that $D^{-1}T$ is unit lower triangular. By item 2 of \Cref{thm:elem-inverse}, every unit lower triangular matrix can be written as $L_1^{-1} \cdots L_{m-1}^{-1}$ (take its strictly lower entries as the multipliers), so its inverse $L_{m-1} \cdots L_1$ is lower triangular. Therefore $T^{-1} = (D^{-1}T)^{-1}D^{-1}$ is a product of lower triangular matrices and is lower triangular.

\paragraph{A shorter proof of item 2 of \Cref{thm:elem-inverse}~\cite[Lecture 20]{TB}.} For $k < j$ we have $\vec{e}_k^{\,\ast}\vec{\ell}_j = 0$, so
\begin{equation}
  (I + \vec{\ell}_k\vec{e}_k^{\,\ast})(I + \vec{\ell}_j\vec{e}_j^{\,\ast}) = I + \vec{\ell}_k\vec{e}_k^{\,\ast} + \vec{\ell}_j\vec{e}_j^{\,\ast},
\end{equation}
and induction gives $L_1^{-1} \cdots L_{m-1}^{-1} = I + \sum_{k=1}^{m-1}\vec{\ell}_k\vec{e}_k^{\,\ast}$.

\paragraph{Cost, and solving with LU.} Step $k$ of \Cref{alg:lu} updates an $(m - k) \times (m - k)$ block with one multiplication and one subtraction per entry, so the total is
\begin{equation}
  \sum_{k=1}^{m-1} 2(m - k)^2 \approx \frac{2}{3}m^3 \text{ flops.}
\end{equation}
Solving $A\vec{x} = \vec{b}$ splits into three steps: factor $A = LU$ ($\tfrac{2}{3}m^3$), solve $L\vec{y} = \vec{b}$ by forward substitution ($m^2$), and solve $U\vec{x} = \vec{y}$ by back substitution ($m^2$). The factorization is done once, and each additional right-hand side costs only $2m^2$. This is the reason not to form $A^{-1}$. The two inner loops form a rank-one update, a Level 2 operation; see \Cref{sec:blocked-lu} for the blocked version.

\paragraph{The pivots as ratios of minors~\cite[\S2.3]{Demmel}.} Taking determinants in the proof of \Cref{thm:lu-exists} gives
\begin{equation}
  U_{kk} = \frac{\det A[1{:}k, 1{:}k]}{\det A[1{:}k{-}1, 1{:}k{-}1]},
\end{equation}
and $A = LU$ with $L$ unit lower triangular exists and is unique if and only if $A[1{:}k, 1{:}k]$ is nonsingular for every $k = 1, \dots, m - 1$.

\paragraph{Pivoting, preview.} In the $\eps$ example, swapping the two rows first gives $\begin{bmatrix} 1 & 1 \\ \eps & 1 \end{bmatrix}$, whose only multiplier is $\eps$, and the problem disappears. This is the idea of \emph{pivoting}: at each step, move the entry of largest absolute value in the current column to the pivot position, obtaining $PA = LU$ with $\abs{\ell_{ik}} \le 1$, which controls $\abs{L}$ in the error bound. It is the subject of \Cref{ch:pivoting}.

\chapter{Error Analysis of LU and Pivoting}
\label{ch:pivoting}

Supplementary reading:
\begin{itemize}
  \item \cite{TB} Lectures 21 and 22.
  \item \cite{Demmel} Sections 2.3 and 2.4.
  \item \cite{Higham} Chapter 9.
\end{itemize}

\section{Backward Error of the LU Factorization}

Recall from \Cref{ch:lu} that Gaussian elimination computes $A = LU$ in place, with the multipliers stored below the diagonal and $U$ on and above it, and that it fails on a zero pivot even when $A$ is nonsingular (for example $\begin{bmatrix} 0 & 1 \\ 1 & 1 \end{bmatrix}$). \Cref{ch:lu} ended with an example in which a tiny pivot destroys the LU factorization of the well-conditioned matrix
\begin{equation}
  A = \begin{bmatrix} \eps & 1 \\ 1 & 1 \end{bmatrix}, \qquad \kappa(A) \approx 2.6 \text{ as } \eps \to 0.
\end{equation}
Here $\eps$ is just a small number, smaller than $\epsm$; it is not the machine epsilon. The problem is well-conditioned, so the failure lies with the algorithm. The following theorem explains exactly where it comes from.

\begin{theorem}[Backward Error of LU]
  \label{thm:lu-backward}
  Suppose the LU factorization of $A \in \K^{m \times m}$ runs to completion in floating-point arithmetic without encountering a zero pivot, producing $\tilde{L}$ and $\tilde{U}$. Then, to first order in $\epsm$,
  \begin{equation}
    A = \tilde{L}\tilde{U} + E, \qquad \abs{E} \le m\epsm\abs{\tilde{L}}\abs{\tilde{U}}.
  \end{equation}
\end{theorem}

In other words, the computed $\tilde{L}$ and $\tilde{U}$ are the exact LU factors of the perturbed matrix $A - E$, and the perturbation is controlled by $\abs{\tilde{L}}\abs{\tilde{U}}$. The analysis ignores terms of order $O(\epsm^2)$.

\subsection{Entrywise Formulas}

From $A = LU$ and $L_{ii} = 1$ we have $A_{ik} = \sum_{j=1}^{\min(i,k)} L_{ij}U_{jk}$. Solving for the entries of $L$ and $U$ gives formulas that produce the same numbers as Gaussian elimination, only in a different order of operations:
\begin{align}
  U_{ik} &= A_{ik} - \sum_{j=1}^{i-1} L_{ij}U_{jk}, & i &\le k, \label{eq:lu-u} \\
  L_{ik} &= \frac{1}{U_{kk}}\left(A_{ik} - \sum_{j=1}^{k-1} L_{ij}U_{jk}\right), & i &> k. \label{eq:lu-l}
\end{align}

\subsection{Proof of \Cref{thm:lu-backward}}

We use the following fact about inner products (see Homework 2 and \Cref{prop:dot}): an inner product of length $m$ computed in floating point satisfies
\begin{equation}
  \fl\left(\sum_{i=1}^{m} x_iy_i\right) = \sum_{i=1}^{m} x_iy_i(1 + \delta_i), \qquad \abs{\delta_i} \lesssim m\epsm,
\end{equation}
so each term carries its own relative perturbation, proportional to the number of terms summed.

\begin{proof}
  \emph{The $U$ part ($i \le k$).} Computing \eqref{eq:lu-u} in floating point gives
  \begin{equation}
    \tilde{U}_{ik} = \left(A_{ik} - \sum_{j=1}^{i-1}\tilde{L}_{ij}\tilde{U}_{jk}(1 + \delta_j)\right)(1 + \delta'), \qquad \abs{\delta_j} \le (i - 1)\epsm, \quad \abs{\delta'} \le \epsm,
  \end{equation}
  where $\delta'$ comes from the final subtraction. Solving for $A_{ik}$ and using $\tilde{L}_{ii} = 1$,
  \begin{equation}
    A_{ik} = \frac{1}{1 + \delta'}\tilde{L}_{ii}\tilde{U}_{ik} + \sum_{j=1}^{i-1}\tilde{L}_{ij}\tilde{U}_{jk}(1 + \delta_j).
  \end{equation}
  Setting $1 + \delta_i \doteq (1 + \delta')^{-1}$, so that $\abs{\delta_i} \lesssim \epsm$ to first order, the $i$-th term joins the sum:
  \begin{equation}
    A_{ik} = \sum_{j=1}^{i}\tilde{L}_{ij}\tilde{U}_{jk} + \underbrace{\sum_{j=1}^{i}\tilde{L}_{ij}\tilde{U}_{jk}\delta_j}_{\doteq\, E_{ik}}.
  \end{equation}
  All $\abs{\delta_j} \le m\epsm$, so
  \begin{equation}
    \abs{E_{ik}} \le \sum_{j=1}^{i}\abs{\tilde{L}_{ij}}\abs{\tilde{U}_{jk}}\,m\epsm = m\epsm\left(\abs{\tilde{L}}\abs{\tilde{U}}\right)_{ik},
  \end{equation}
  where the last step uses that $\tilde{L}$ is lower triangular, so the terms with $j > i$ vanish and summing to $i$ is the same as summing to $m$.

  \emph{The $L$ part ($i > k$).} Similarly, computing \eqref{eq:lu-l} in floating point gives
  \begin{equation}
    \tilde{L}_{ik} = (1 + \delta'')\,\frac{(1 + \delta')\left(A_{ik} - \sum_{j=1}^{k-1}\tilde{L}_{ij}\tilde{U}_{jk}(1 + \delta_j)\right)}{\tilde{U}_{kk}}, \qquad \abs{\delta''} \le \epsm,
  \end{equation}
  with $\delta'$ from the subtraction and $\delta''$ from the division. Solving for $A_{ik}$ and setting $1 + \delta_k \doteq \left((1 + \delta')(1 + \delta'')\right)^{-1}$, so that $\abs{\delta_k} \lesssim 2\epsm \le m\epsm$, we get
  \begin{equation}
    A_{ik} = \sum_{j=1}^{k}\tilde{L}_{ij}\tilde{U}_{jk} + E_{ik}, \qquad \abs{E_{ik}} \le m\epsm\left(\abs{\tilde{L}}\abs{\tilde{U}}\right)_{ik}.
  \end{equation}
  The two cases together prove the theorem.
\end{proof}

\subsection{What the Bound Says}

If $\abs{\tilde{L}}\abs{\tilde{U}}$ is of the same order as $\abs{A}$, then $E$ is a relative perturbation of size $O(m\epsm)$ and the algorithm is backward stable. But $\abs{\tilde{L}}\abs{\tilde{U}}$ can be far larger than $\abs{A}$. In the example above,
\begin{equation}
  \abs{L}\abs{U} = \begin{bmatrix} 1 & 0 \\ \eps^{-1} & 1 \end{bmatrix}\begin{bmatrix} \eps & 1 \\ 0 & \eps^{-1} - 1 \end{bmatrix} = \begin{bmatrix} \eps & 1 \\ 1 & 2\eps^{-1} - 1 \end{bmatrix}.
\end{equation}
The bottom-right entry is about $2\eps^{-1}$; multiplied by $\epsm$, the error can be $O(1)$ or larger, exactly as observed. The question becomes: \emph{how can we keep $\abs{L}\abs{U}$ small?}

\section{Pivots and the Growth of Products}

Look at the storage after $k - 1$ steps of elimination: the first $k - 1$ rows already hold $U$, the first $k - 1$ columns below the diagonal hold $L$, and the trailing block $\tilde{A}$ is still being processed.
\begin{equation}
  \begin{bmatrix}
    U_{11} & \cdots & & & \\
    & \ddots & & & \\
    & & U_{k-1,k-1} & U_{k-1,k} \ \cdots & U_{k-1,m} \\
    & L & L_{k,k-1} & \tilde{A}_{kk} \ \cdots & \tilde{A}_{km} \\
    & & \vdots & \vdots & \vdots \\
    & & L_{m,k-1} & \tilde{A}_{mk} \ \cdots & \tilde{A}_{mm}
  \end{bmatrix}
\end{equation}
Step $k$ turns row $k$ into a row of $U$ and column $k$ below the diagonal into a column of $L$, and updates the trailing $(m - k) \times (m - k)$ block:
\begin{equation}
  L_{ik} = \frac{\tilde{A}_{ik}}{\tilde{A}_{kk}}, \qquad U_{kj} = \tilde{A}_{kj}, \qquad \tilde{A}_{ij} \gets \tilde{A}_{ij} - L_{ik}U_{kj}.
\end{equation}
The entry $\tilde{A}_{kk}$ is the $k$-th \emph{pivot}. The largest product created in step $k$, which enters both the update and $\abs{L}\abs{U}$, is
\begin{equation}
  \max_{i,j}\abs{L_{ik}U_{kj}} = \max_{i,j}\frac{\abs{\tilde{A}_{ik}}\abs{\tilde{A}_{kj}}}{\abs{\tilde{A}_{kk}}}.
\end{equation}
The numerator is determined by column $k$ and row $k$. To keep this quantity small, we want the denominator, the pivot, to be as large as possible.

\section{Partial Pivoting}

We must pivot when $A_{kk} = 0$, since otherwise the algorithm cannot continue. Even when $A_{kk} \ne 0$, we want to avoid \emph{small} pivots, which make the products $L_{ik}U_{kj}$ explode.

\begin{definition}[Partial (Row) Pivoting]
  At step $k$ of Gaussian elimination, \emph{partial pivoting} chooses as pivot the entry of largest absolute value in column $k$, on or below the diagonal:
  \begin{equation}
    r = \operatorname*{argmax}_{i \ge k}\abs{\tilde{A}_{ik}}.
  \end{equation}
  It swaps rows $r$ and $k$ with a permutation matrix $P_k$ and then eliminates with $L_k$. The sequence $A \xrightarrow{L_k} \cdots$ becomes $A \xrightarrow{P_k} \cdot \xrightarrow{L_k} \cdots$.
\end{definition}

Because the pivot is the largest entry in its column, every multiplier satisfies
\begin{equation}
  \abs{L_{ik}} = \frac{\abs{\tilde{A}_{ik}}}{\abs{\tilde{A}_{kk}}} \le 1,
\end{equation}
so $\abs{L_{ik}U_{kj}} \le \abs{U_{kj}}$: no update exceeds the size of the current pivot row.

\subsection{Moving the Permutations Together}

With row swaps, elimination becomes
\begin{equation}
  L_{m-1}P_{m-1} \cdots L_2P_2L_1P_1A = U,
\end{equation}
where $L$'s and $P$'s alternate, so we cannot directly write $PA = LU$. We need to move all permutations to the right.

\begin{proposition}[Commuting a Permutation Past an Elimination]
  \label{prop:perm-swap}
  Let $\ell > k$ and let $P_\ell$ swap rows $\ell$ and $i$ for some $i \ge \ell$, so that $P_\ell = P_\ell^{-1}$. Then
  \begin{equation}
    P_\ell L_k = (P_\ell L_kP_\ell)P_\ell, \qquad P_\ell L_kP_\ell = I - (P_\ell\vec{\ell}_k)\vec{e}_k^{\,\ast},
  \end{equation}
  so $P_\ell L_kP_\ell$ is again an elementary lower triangular matrix, obtained from $L_k$ by swapping two of its multipliers.
\end{proposition}

\begin{proof}
  The first identity is $P_\ell^2 = I$. Since $\ell, i > k$, the permutation $P_\ell$ fixes $\vec{e}_k$, so $\vec{e}_k^{\,\ast}P_\ell = \vec{e}_k^{\,\ast}$ and
  \begin{equation}
    P_\ell(I - \vec{\ell}_k\vec{e}_k^{\,\ast})P_\ell = I - (P_\ell\vec{\ell}_k)(\vec{e}_k^{\,\ast}P_\ell) = I - (P_\ell\vec{\ell}_k)\vec{e}_k^{\,\ast}.
  \end{equation}
  Geometrically, left multiplication by $P_\ell$ swaps two rows of $L_k$, moving two diagonal ones off the diagonal, and right multiplication swaps the corresponding columns, moving them back; the net effect is to swap the multipliers in rows $\ell$ and $i$ of column $k$.
\end{proof}

Using \Cref{prop:perm-swap} repeatedly to push every $P$ to the right,
\begin{equation}
  L_{m-1}P_{m-1} \cdots L_1P_1 = L'_{m-1} \cdots L'_1\,P_{m-1} \cdots P_1, \qquad L'_k \doteq P_{m-1} \cdots P_{k+1}L_kP_{k+1} \cdots P_{m-1},
\end{equation}
and each $L'_k$ is still elementary lower triangular. With $P \doteq P_{m-1} \cdots P_1$ and $L \doteq (L'_{m-1} \cdots L'_1)^{-1}$, we obtain
\begin{equation}
  PA = LU.
\end{equation}

\section{The In-Place Algorithm with Partial Pivoting}

How do we compute this in place, up to the permutation? Suppose $k$ steps are done:
\begin{equation}
  L_k \cdots L_1PA = U_k,
\end{equation}
where the $L_j$ are already the ``permutation-adjusted'' versions and $U_k$ is the current working matrix. In storage, $U_k$ occupies the upper triangle and the trailing block, and the multipliers of $L_1, \dots, L_k$ sit in the strictly lower part of the first $k$ columns. Step $k + 1$ swaps rows and then eliminates:
\begin{equation}
  L_{k+1}P_{k+1}L_k \cdots L_1PA = L_{k+1}P_{k+1}U_k.
\end{equation}
Moving $P_{k+1}$ to the right,
\begin{equation}
  L_{k+1}\underbrace{L'_k \cdots L'_1}_{\text{multipliers permuted by } P_{k+1}}\underbrace{P_{k+1}P}_{P'}A = L_{k+1}\underbrace{P_{k+1}U_k}_{U'_k} = U_{k+1}.
\end{equation}
Reading term by term:
\begin{itemize}
  \item $L'_j = P_{k+1}L_jP_{k+1}$ only permutes the stored multipliers;
  \item $P' = P_{k+1}P$ swaps the same two rows of the accumulated permutation;
  \item $U'_k = P_{k+1}U_k$ swaps the same two rows of the working matrix.
\end{itemize}
Since the multipliers and $U_k$ live in the same array, \emph{swapping entire rows (columns $1$ through $m$)} performs all three updates at once, after which we eliminate as usual.

\begin{algorithm}[H]
  \caption{LU factorization with partial pivoting, in place.}
  \label{alg:gepp}
  \begin{algorithmic}[1]
    \State $P \gets I$
    \For{$k = 1, \dots, m - 1$}
      \State $r \gets \operatorname{argmax}_{i \ge k}\abs{A_{ik}}$
      \State swap rows $r$ and $k$ of $A$ (all columns $1, \dots, m$) and of $P$
      \For{$i = k + 1, \dots, m$}
        \State $A_{ik} \gets A_{ik}/A_{kk}$
      \EndFor
      \For{$j = k + 1, \dots, m$}
        \For{$i = k + 1, \dots, m$}
          \State $A_{ij} \gets A_{ij} - A_{ik}A_{kj}$
        \EndFor
      \EndFor
    \EndFor
  \end{algorithmic}
\end{algorithm}

Some remarks:
\begin{itemize}
  \item In Julia, \texttt{argmax(abs.(A[k:end, k]))} returns an index into the subvector, so one adds $k - 1$ to recover the global row index.
  \item The row swap runs over all columns $1, \dots, m$, not just $k, \dots, m$: the first $k - 1$ columns hold multipliers, and swapping them is exactly $L'_j = P_{k+1}L_jP_{k+1}$.
  \item On exit the strictly lower part of $A$ holds $L$, the upper part holds $U$, and $PA = LU$.
  \item In practice $P$ is not stored as a matrix but as a permutation vector of length $m$ (LAPACK's \texttt{ipiv}), so each swap exchanges two integers. The search for pivots costs only $O(m^2)$ comparisons on top of the $\tfrac{2}{3}m^3$ flops, although in column-major storage a row swap is a strided access whose cost is not negligible.
\end{itemize}

\section{(SUPPLEMENT) Growth Factor}

\begin{definition}[Growth Factor]
  The \emph{growth factor} of Gaussian elimination on $A$ is
  \begin{equation}
    \rho \doteq \frac{\max_{i,j}\abs{U_{ij}}}{\max_{i,j}\abs{A_{ij}}}.
  \end{equation}
\end{definition}

With partial pivoting $\abs{L_{ij}} \le 1$, so $\abs{L}\abs{U}$ is governed by $U$, and the backward error is roughly $\norm{E} \lesssim m^2\rho\,\epsm\norm{A}$ (the constant depends on the norm). Since $\abs{L_{ik}} \le 1$ in each update $\tilde{A}_{ij} - L_{ik}U_{kj}$, each step can at most double an entry, so $\rho \le 2^{m-1}$. This worst case is attained~\cite[Lecture 22]{TB}; \Cref{ch:ls} opens with the example. Yet on matrices met in practice $\rho$ is almost always small, which is why partial pivoting is used as a stable algorithm in practice.

\section{(SUPPLEMENT) Complete Pivoting}

Complete pivoting chooses the entry of largest absolute value in the entire trailing block as the pivot and swaps both rows and columns, giving $PAQ = LU$. Its growth factor obeys a much better (subexponential) bound, but each step searches $O((m - k)^2)$ entries, adding $O(m^3)$ comparisons in total, so it is rarely used.

\chapter{Least Squares and Gram--Schmidt Orthogonalization}
\label{ch:ls}

Supplementary reading:
\begin{itemize}
  \item \cite{TB} Lectures 6, 7, 8, 11, and 22.
  \item \cite{Demmel} Sections 3.1--3.3.
\end{itemize}

\section{Is Partial Pivoting Enough?}

\Cref{ch:pivoting} bounded the backward error of LU by $m\epsm\abs{\tilde{L}}\abs{\tilde{U}}$ and introduced partial pivoting to keep the multipliers bounded, $\abs{L_{ij}} \le 1$; commuting the permutations past the elimination matrices then gave $PA = LU$ and an in-place algorithm. Rows and columns can also be pivoted simultaneously (\emph{complete pivoting}), giving $PA\tilde{P} = LU$ with a column permutation $\tilde{P}$. In practice, partial pivoting is usually enough. In theory, however, it can be \emph{unstable} when $m$ is large.

\begin{example}[Exponential Growth Under Partial Pivoting]
  \label{ex:growth}
  Consider
  \begin{equation}
    A = \begin{bmatrix} 1 & & & 1 \\ -1 & 1 & & 1 \\ -1 & -1 & 1 & 1 \\ -1 & -1 & -1 & 1 \end{bmatrix}.
  \end{equation}
  In every column the candidate pivots have the same absolute value, $\abs{1} = \abs{-1}$, and \texttt{argmax} returns the first, so partial pivoting never swaps rows. Tracking the in-place storage (multipliers below the diagonal),
  \begin{equation}
    A \rightsquigarrow
    \begin{bmatrix} 1 & & & 1 \\ -1 & 1 & & 2 \\ -1 & -1 & 1 & 2 \\ -1 & -1 & -1 & 2 \end{bmatrix} \rightsquigarrow
    \begin{bmatrix} 1 & & & 1 \\ -1 & 1 & & 2 \\ -1 & -1 & 1 & 4 \\ -1 & -1 & -1 & 4 \end{bmatrix} \rightsquigarrow
    \begin{bmatrix} 1 & & & 1 \\ -1 & 1 & & 2 \\ -1 & -1 & 1 & 4 \\ -1 & -1 & -1 & 8 \end{bmatrix}.
  \end{equation}
  The multiplier at step $k$ is $-1$, i.e., row $k$ is added to every row below it. Row $k$ is zero in columns $k + 1$ through $m - 1$, so the middle part never changes, but the last column doubles at every step. The result is
  \begin{equation}
    L = \begin{bmatrix} 1 & & & \\ -1 & 1 & & \\ -1 & -1 & 1 & \\ -1 & -1 & -1 & 1 \end{bmatrix}, \qquad
    U = \begin{bmatrix} 1 & & & 1 \\ & 1 & & 2 \\ & & 1 & 4 \\ & & & 8 \end{bmatrix}.
  \end{equation}
  For general $m$, $U_{mm} = 2^{m-1}$, so $\abs{L}\abs{U} \gtrsim 2^m$ and the error bound $m\epsm\abs{L}\abs{U}$ can reach $m2^m\epsm$. In double precision, all accuracy is lost for $m \approx 60$.
\end{example}

This is the worst case $\rho = 2^{m-1}$ of the growth factor mentioned in \Cref{ch:pivoting}. Fortunately, such matrices are extremely rare in practice.

\section{Non-Square Systems}

So far we have only considered $A \in \K^{m \times n}$ with $m = n$. If $m \ne n$, there are two cases.

\paragraph{Overdetermined systems ($m > n$).} $A$ is tall: there are more equations than unknowns. Then $\range(A) \ne \K^m$, since $n$ columns span at most an $n$-dimensional subspace. Hence for a general $\vec{b}$ the equation $A\vec{x} = \vec{b}$ \emph{need not have a solution}. Such a system is called an \emph{overdetermined linear system}.

\paragraph{Underdetermined systems ($m < n$).} $A$ is wide: there are more unknowns than equations. Then $\nullsp(A) \ne \{\zerovec\}$, since the rank is at most $m < n$. Hence even when a solution exists, it is \emph{never unique}: adding any vector of the null space to a solution gives another solution. Such a system is called an \emph{underdetermined linear system}.

\section{Least Squares Problems}

Each case lacks one property, existence or uniqueness, and a minimization problem restores it:
\begin{equation}
  \text{overdetermined:}\quad \min_{\vec{x}} \norm{A\vec{x} - \vec{b}}, \qquad\qquad \text{underdetermined:}\quad \min_{\vec{x} \,:\, A\vec{x} = \vec{b}} \norm{\vec{x}}.
\end{equation}
In the overdetermined case we look for the $\vec{x}$ that comes closest to satisfying the equations; in the underdetermined case we pick the ``smallest'' of all solutions.

\subsection{Which Norm?}

There are many choices, depending on the application. In linear regression, for instance, the norm in the overdetermined case says how we measure the \emph{fitting error}, and in the underdetermined case how we measure the \emph{complexity of the model}.

The key fact is that \emph{the solution depends linearly on $\vec{b}$ only when the norm is induced by an inner product}, $\norm{\cdot} = \sqrt{\ip{\cdot}{\cdot}}$. This gives the \emph{least squares problem}, which can be solved with the tools of linear algebra.

\begin{example}[Mean Versus Median]
  Let $A = \begin{bmatrix} 1 & 1 & 1 \end{bmatrix}^\T$, so that $\min_x \norm{Ax - \vec{b}}$ fits a single number to $b_1, b_2, b_3$. In the $2$-norm the answer is the mean $\frac{1}{3}(b_1 + b_2 + b_3)$, which is linear in $\vec{b}$. In the $1$-norm the answer is the median, which is not a linear function of $\vec{b}$.
\end{example}

\section{The Geometry of Least Squares}

For $m = n$, solving $A\vec{x} = \vec{b}$ means finding a linear combination of the columns of $A$ that equals $\vec{b}$ exactly:
\begin{equation}
  A\vec{x} = x_1\vec{a}_1 + x_2\vec{a}_2 + \cdots + x_n\vec{a}_n = \vec{b}.
\end{equation}
For $m > n$, consider
\begin{equation}
  \min_{\vec{x}} \norm{A\vec{x} - \vec{b}}_2^2.
\end{equation}
The columns of $A$ span only an $n$-dimensional subspace $\range(A)$ of $\K^m$, and $\vec{b}$ generally does not lie in it. As $\vec{x}$ varies, $A\vec{x}$ ranges over this subspace, and the point closest to $\vec{b}$ is the \emph{orthogonal projection} of $\vec{b}$ onto $\range(A)$: dropping a perpendicular from $\vec{b}$ to the subspace, its foot is the optimal $A\vec{x}$, and the residual $A\vec{x} - \vec{b}$ is orthogonal to the whole subspace.

\begin{center}
\begin{tikzpicture}[scale=1.1]
  \fill[black!6] (0,0) -- (5,0) -- (6.2,1.6) -- (1.2,1.6) -- cycle;
  \node at (5.3,0.3) {$\range(A)$};
  \coordinate (O) at (1.2,0.4);
  \coordinate (P) at (4,0.9);
  \coordinate (B) at (4,3);
  \draw[-{Latex},thick] (O) -- (B) node[above left] {$\vec{b}$};
  \draw[-{Latex},thick] (O) -- (P) node[below right] {$A\vec{x}$};
  \draw[dashed] (P) -- (B) node[midway,right] {$\vec{b} - A\vec{x}$};
  \draw (4,1.15) -- (3.78,1.1) -- (3.78,0.86);
\end{tikzpicture}
\end{center}

Writing ``the residual is orthogonal to every column of $A$'' as an equation gives the \emph{normal equations}:
\begin{equation}
  A^\ast(A\vec{x} - \vec{b}) = \zerovec \quad\iff\quad A^\ast A\vec{x} = A^\ast\vec{b}. \label{eq:normal}
\end{equation}

\subsection{Orthonormal Columns}

If the columns of $A$ are orthonormal ($A^\ast A = I$), the solution is particularly simple:
\begin{equation}
  \vec{x} = A^\ast\vec{b}.
\end{equation}
Intuitively, $x_i = \ip{\vec{a}_i}{\vec{b}}$ measures the ``length'' of $\vec{b}$ in the ``direction'' of the $i$-th column. Since the columns are mutually orthogonal, the component along each direction can be computed \emph{independently of the others} and the results simply added.

\begin{example}[Projection onto a Coordinate Plane]
  For
  \begin{equation}
    A = \begin{bmatrix} 1 & 0 \\ 0 & 1 \\ 0 & 0 \end{bmatrix},
  \end{equation}
  we get $A^\ast\vec{b} = (b_1, b_2)^\T$, the coordinates of the projection of $\vec{b}$ onto the $xy$-plane; the discarded component $b_3$ is the residual.
\end{example}

\section{Orthogonalization: The QR Factorization}

What if the columns of $A$ are not orthonormal? The idea is to find an \emph{orthonormal basis} of the space spanned by the columns of $A$ and to transfer the problem to that basis. In matrix form,
\begin{equation}
  A = QR,
\end{equation}
where $Q \in \K^{m \times n}$ has orthonormal columns and $R \in \K^{n \times n}$ is invertible (it will turn out to be upper triangular). The least squares solution is then
\begin{equation}
  \vec{x} = R^{-1}Q^\ast\vec{b}.
\end{equation}
Both steps have a clear meaning:
\begin{itemize}
  \item $Q^\ast\vec{b}$ expresses $\vec{b}$ in terms of the columns of $Q$ (the coordinates of the projection, exactly as in the orthonormal case);
  \item $A = QR$ says that $R$ expresses the columns of $A$ as combinations of the columns of $Q$, so $R^{-1}$ converts coefficients with respect to $Q$ back into coefficients with respect to $A$.
\end{itemize}
Algebraically, in the normal equations \eqref{eq:normal} we have $A^\ast A = R^\ast Q^\ast QR = R^\ast R$ and $A^\ast\vec{b} = R^\ast Q^\ast\vec{b}$; cancelling the invertible $R^\ast$ leaves $R\vec{x} = Q^\ast\vec{b}$.

\section{Classical Gram--Schmidt}

\begin{algorithm}[H]
  \caption{Classical Gram--Schmidt (CGS).}
  \label{alg:cgs}
  \begin{algorithmic}[1]
    \For{$j = 1, \dots, n$}
      \State $\hat{\vec{q}}_j \gets \vec{a}_j$
      \For{$i = 1, \dots, j - 1$}
        \State $r_{ij} \gets \ip{\vec{q}_i}{\vec{a}_j}$ \Comment{inner product with the \emph{original} column}
        \State $\hat{\vec{q}}_j \gets \hat{\vec{q}}_j - r_{ij}\vec{q}_i$
      \EndFor
      \State $r_{jj} \gets \norm{\hat{\vec{q}}_j}_2$
      \State $\vec{q}_j \gets \hat{\vec{q}}_j / r_{jj}$
    \EndFor
  \end{algorithmic}
\end{algorithm}

From $\vec{a}_j$ we subtract its projections onto the already computed $\vec{q}_1, \dots, \vec{q}_{j-1}$ and normalize what is left. The coefficients $r_{ij}$ form an upper triangular matrix $R$, because
\begin{equation}
  \vec{a}_j = r_{jj}\vec{q}_j + \sum_{i=1}^{j-1} r_{ij}\vec{q}_i.
\end{equation}

\begin{theorem}[Gram--Schmidt]
  \label{thm:gs}
  Let $\vec{a}_1, \dots, \vec{a}_n \in \K^m$ ($m \ge n$) be linearly independent. Then the vectors $\vec{q}_1, \dots, \vec{q}_n$ produced by \Cref{alg:cgs} form an orthonormal basis of $\spn(\vec{a}_1, \dots, \vec{a}_n)$, and $\spn(\vec{q}_1, \dots, \vec{q}_j) = \spn(\vec{a}_1, \dots, \vec{a}_j)$ for every $j$.
\end{theorem}

\begin{proof}
  By induction. By construction, $\hat{\vec{q}}_j, \vec{q}_j \in \spn(\vec{a}_1, \dots, \vec{a}_j)$. Suppose $\vec{q}_1, \dots, \vec{q}_j$ are orthonormal. For every $k \le j$,
  \begin{equation}
    \ip{\vec{q}_k}{\vec{a}_{j+1} - \sum_{i=1}^{j}\ip{\vec{q}_i}{\vec{a}_{j+1}}\vec{q}_i} = \ip{\vec{q}_k}{\vec{a}_{j+1}} - \ip{\vec{q}_k}{\vec{a}_{j+1}} = 0,
  \end{equation}
  using $\ip{\vec{q}_k}{\vec{q}_i} = \delta_{ki}$, so that only the term $i = k$ survives in the sum. Hence $\hat{\vec{q}}_{j+1}$ is orthogonal to $\vec{q}_1, \dots, \vec{q}_j$. Moreover $\vec{a}_{j+1} \notin \spn(\vec{a}_1, \dots, \vec{a}_j) = \spn(\vec{q}_1, \dots, \vec{q}_j)$, so $\hat{\vec{q}}_{j+1} \ne \zerovec$ and can be normalized. Thus $\vec{q}_1, \dots, \vec{q}_{j+1}$ are orthonormal. They are $j + 1$ linearly independent vectors in the $(j + 1)$-dimensional space $\spn(\vec{a}_1, \dots, \vec{a}_{j+1})$, so they form a basis of it.
\end{proof}

\begin{example}[CGS Loses Orthogonality]
  \label{ex:cgs-lauchli}
  Consider the matrix
  \begin{equation}
    A = \begin{bmatrix} 1 & 1 & 1 \\ \eps & 0 & 0 \\ 0 & \eps & 0 \\ 0 & 0 & \eps \end{bmatrix},
  \end{equation}
  with $\eps$ small enough that $1 + \eps^2 = 1$ in floating point, i.e., $\eps < \sqrt{\epsm}$, for example $\eps = 10^{-10}$.

  \emph{Step 1.} $\hat{\vec{q}}_1 = (1, \eps, 0, 0)^\T$ and $\vec{q}_1 = (1 + \eps^2)^{-1/2}(1, \eps, 0, 0)^\T$, which in floating point is $(1, \eps, 0, 0)^\T$.

  \emph{Step 2.} $r_{12} = \ip{\vec{q}_1}{\vec{a}_2} = 1$, so
  \begin{equation}
    \hat{\vec{q}}_2 = \begin{bmatrix} 1 \\ 0 \\ \eps \\ 0 \end{bmatrix} - \begin{bmatrix} 1 \\ \eps \\ 0 \\ 0 \end{bmatrix} = \begin{bmatrix} 0 \\ -\eps \\ \eps \\ 0 \end{bmatrix}, \qquad
    \vec{q}_2 = \frac{1}{\sqrt{2}\,\eps}\hat{\vec{q}}_2 = \begin{bmatrix} 0 \\ -1/\sqrt{2} \\ 1/\sqrt{2} \\ 0 \end{bmatrix}.
  \end{equation}

  \emph{Step 3.} $r_{13} = \ip{\vec{q}_1}{\vec{a}_3} = 1$ and $r_{23} = \ip{\vec{q}_2}{\vec{a}_3} = 0$, so
  \begin{equation}
    \hat{\vec{q}}_3 = \begin{bmatrix} 1 \\ 0 \\ 0 \\ \eps \end{bmatrix} - \begin{bmatrix} 1 \\ \eps \\ 0 \\ 0 \end{bmatrix} - 0 \cdot \vec{q}_2 = \begin{bmatrix} 0 \\ -\eps \\ 0 \\ \eps \end{bmatrix}, \qquad
    \vec{q}_3 = \begin{bmatrix} 0 \\ -1/\sqrt{2} \\ 0 \\ 1/\sqrt{2} \end{bmatrix}.
  \end{equation}
  The result
  \begin{equation}
    Q = \begin{bmatrix} 1 & 0 & 0 \\ \eps & -1/\sqrt{2} & -1/\sqrt{2} \\ 0 & 1/\sqrt{2} & 0 \\ 0 & 0 & 1/\sqrt{2} \end{bmatrix}
  \end{equation}
  is \emph{far from orthogonal}: $\ip{\vec{q}_2}{\vec{q}_3} = \frac{1}{2}$, so the angle between them is $60^\circ$ rather than $90^\circ$.
\end{example}

\section{Modified Gram--Schmidt}

\paragraph{What went wrong.} Subtracting $\vec{q}_1$ from $\vec{a}_3$ subtracts two $O(1)$ vectors and leaves only an $O(\eps)$ remainder: severe cancellation, so the \emph{relative error} of the result is large. More precisely, the computed $\vec{q}_1$ and $\vec{q}_2$ are not exactly orthogonal: $\ip{\vec{q}_1}{\vec{q}_2} = -\eps/\sqrt{2}$. This is tiny in absolute terms, but $\vec{a}_3 - \vec{q}_1$ is itself only $O(\eps)$, and relative to it the deviation is $O(1)$.

CGS computes $r_{23}$ as the inner product of $\vec{q}_2$ with the \emph{original} $\vec{a}_3$. In exact arithmetic $\ip{\vec{q}_2}{\vec{a}_3} = \ip{\vec{q}_2}{\vec{a}_3 - r_{13}\vec{q}_1}$, because $\vec{q}_2 \perp \vec{q}_1$, so there is no difference; in floating point this identity no longer holds. What we actually want is
\begin{equation}
  \vec{q}_2 \perp \vec{a}_3 - \vec{q}_1 \qquad \text{rather than} \qquad \vec{q}_2 \perp \vec{a}_3.
\end{equation}

\paragraph{The fix.} Change a single line: compute $r_{ij}$ from the \emph{partially orthogonalized} $\hat{\vec{q}}_j$ instead of the original $\vec{a}_j$.

\begin{algorithm}[H]
  \caption{Modified Gram--Schmidt (MGS).}
  \label{alg:mgs}
  \begin{algorithmic}[1]
    \For{$j = 1, \dots, n$}
      \State $\hat{\vec{q}}_j \gets \vec{a}_j$
      \For{$i = 1, \dots, j - 1$}
        \State $r_{ij} \gets \ip{\vec{q}_i}{\hat{\vec{q}}_j}$ \Comment{inner product with the \emph{updated} vector}
        \State $\hat{\vec{q}}_j \gets \hat{\vec{q}}_j - r_{ij}\vec{q}_i$
      \EndFor
      \State $r_{jj} \gets \norm{\hat{\vec{q}}_j}_2$
      \State $\vec{q}_j \gets \hat{\vec{q}}_j / r_{jj}$
    \EndFor
  \end{algorithmic}
\end{algorithm}

In exact arithmetic CGS and MGS are identical. They differ only in floating point: MGS projects out each direction from \emph{what is left after the previous projections}, so errors made in earlier steps are corrected by later ones.

\begin{remark}[Orthogonality Is Fragile]
  Arguments based on orthogonality are often fragile in floating-point arithmetic. The identity $\ip{\vec{q}_i}{\vec{q}_k} = 0$ holds only approximately, and any manipulation that relies on it can fail as soon as cancellation occurs.
\end{remark}

\begin{example}[MGS on the Same Matrix]
  For the matrix of \Cref{ex:cgs-lauchli}, $\vec{q}_1$ and $\vec{q}_2$ are computed exactly as in CGS. Step 3 now projects in two stages:
  \begin{equation}
    \hat{\vec{q}}_3^{(1)} = \begin{bmatrix} 1 \\ 0 \\ 0 \\ \eps \end{bmatrix} - \begin{bmatrix} 1 \\ \eps \\ 0 \\ 0 \end{bmatrix} = \begin{bmatrix} 0 \\ -\eps \\ 0 \\ \eps \end{bmatrix}.
  \end{equation}
  Now $r_{23} = \ip{\vec{q}_2}{\hat{\vec{q}}_3^{(1)}} = \eps/\sqrt{2}$ is no longer zero, and
  \begin{equation}
    \hat{\vec{q}}_3^{(2)} = \begin{bmatrix} 0 \\ -\eps \\ 0 \\ \eps \end{bmatrix} - \frac{\eps}{\sqrt{2}}\begin{bmatrix} 0 \\ -1/\sqrt{2} \\ 1/\sqrt{2} \\ 0 \end{bmatrix} = \begin{bmatrix} 0 \\ -\eps/2 \\ -\eps/2 \\ \eps \end{bmatrix}, \qquad
    \vec{q}_3 = \frac{\hat{\vec{q}}_3^{(2)}}{\eps\sqrt{3/2}} = \begin{bmatrix} 0 \\ -1/\sqrt{6} \\ -1/\sqrt{6} \\ \sqrt{2/3} \end{bmatrix}.
  \end{equation}
  The result is
  \begin{equation}
    Q = \begin{bmatrix} 1 & 0 & 0 \\ \eps & -1/\sqrt{2} & -1/\sqrt{6} \\ 0 & 1/\sqrt{2} & -1/\sqrt{6} \\ 0 & 0 & \sqrt{2/3} \end{bmatrix}.
  \end{equation}
  Now $\ip{\vec{q}_2}{\vec{q}_3} = \frac{1}{\sqrt{12}} - \frac{1}{\sqrt{12}} = 0$, and the remaining inner products $\ip{\vec{q}_1}{\vec{q}_2} = -\eps/\sqrt{2}$ and $\ip{\vec{q}_1}{\vec{q}_3} = -\eps/\sqrt{6}$ are $O(\eps)$: the columns are \emph{nearly orthogonal}.
\end{example}

\section{(SUPPLEMENT) Further Topics}

\paragraph{Gram--Schmidt in a first linear algebra course.} The usual textbook formula, with unnormalized vectors,
\begin{equation}
  \vec{v}_3 = \vec{x}_3 - \frac{\vec{x}_3 \cdot \vec{v}_1}{\vec{v}_1 \cdot \vec{v}_1}\vec{v}_1 - \frac{\vec{x}_3 \cdot \vec{v}_2}{\vec{v}_2 \cdot \vec{v}_2}\vec{v}_2,
\end{equation}
takes both inner products with the \emph{original} $\vec{x}_3$, so it is classical Gram--Schmidt. The modified version proceeds in two stages,
\begin{equation}
  \vec{v}_3^{(1)} = \vec{x}_3 - \frac{\vec{x}_3 \cdot \vec{v}_1}{\vec{v}_1 \cdot \vec{v}_1}\vec{v}_1, \qquad \vec{v}_3 = \vec{v}_3^{(1)} - \frac{\vec{v}_3^{(1)} \cdot \vec{v}_2}{\vec{v}_2 \cdot \vec{v}_2}\vec{v}_2,
\end{equation}
and the two agree in exact arithmetic because $\vec{v}_3^{(1)} \cdot \vec{v}_2 = \vec{x}_3 \cdot \vec{v}_2$ when $\vec{v}_1 \cdot \vec{v}_2 = 0$. A course working in exact arithmetic loses nothing by writing CGS, but the step uses $\vec{v}_1 \cdot \vec{v}_2 = 0$, which is exactly what fails in \Cref{ex:cgs-lauchli}. For $\vec{v}_2$ only one projection is subtracted, so CGS and MGS differ only from the third vector on.

\paragraph{Quantifying the loss of orthogonality.} A result of Bj\"orck states that MGS produces $Q$ with $\norm{I - Q^\ast Q} \approx \epsm\kappa(A)$, while for CGS the loss can be as bad as $\epsm\kappa(A)^2$. The matrix of \Cref{ex:cgs-lauchli} (the L\"auchli matrix) has $\kappa(A) \approx \sqrt{3}/\eps$, so MGS loses about $\epsm/\eps \sim \eps$ when $\eps \approx \sqrt{\epsm}$, in agreement with $\ip{\vec{q}_1}{\vec{q}_2} = O(\eps)$; CGS loses everything.

\paragraph{Why not solve the normal equations?} Solving $A^\ast A\vec{x} = A^\ast\vec{b}$ looks most direct, but $\kappa(A^\ast A) = \kappa(A)^2$~\cite[Lecture 11]{TB}. For the matrix of \Cref{ex:cgs-lauchli},
\begin{equation}
  A^\ast A = \begin{bmatrix} 1 + \eps^2 & 1 & 1 \\ 1 & 1 + \eps^2 & 1 \\ 1 & 1 & 1 + \eps^2 \end{bmatrix},
\end{equation}
and in floating point $1 + \eps^2 = 1$, so the computed $A^\ast A$ is the singular all-ones matrix and all information is lost. Forming $A^\ast A$ discards information that is still present in $A$; QR works on $A$ directly~\cite[\S3.3]{Demmel}.

\paragraph{Householder QR and reorthogonalization.} Even for MGS, the orthogonality of $Q$ depends on $\kappa(A)$. In practice QR is usually computed with Householder reflections (\Cref{ch:householder}), whose $Q$ is orthogonal to working precision independently of $\kappa(A)$. Another common remedy is to run CGS twice (``twice is enough''), which parallelizes better than MGS because the inner products of CGS can be computed at once as a matrix--vector product $Q^\ast\vec{a}_j$.

\chapter{Factorizations as Products and Householder Triangularization}
\label{ch:householder}

Supplementary reading:
\begin{itemize}
  \item \cite{TB} Lectures 8, 10, and 16.
  \item \cite{Demmel} Section 3.4.
\end{itemize}

\section{Why Write Algorithms as Matrix Products?}

The LU factorization and Gram--Schmidt were described in \Cref{ch:lu,ch:ls} as row-by-row or vector-by-vector operations, and the error analysis of \Cref{ch:pivoting} worked entry by entry ($\abs{E} \le m\epsm\abs{L}\abs{U}$). Such an analysis must be redone for every algorithm, and it hardly reveals \emph{why} different algorithms differ in stability.

This lecture takes a different angle: write every factorization algorithm as a \emph{product of a sequence of matrices}, and look only at the properties of the matrix applied at each step. Stability then reduces to one simple quantity, the condition number of these transformation matrices.

\section{The Product Form of Gram--Schmidt}

\subsection{The Analogy with LU}

LU \emph{left}-multiplies by a sequence of lower triangular matrices in order to reach an \emph{upper triangular} matrix:
\begin{equation}
  L_{m-1} \cdots L_1 A = U.
\end{equation}
Each $L_k$ performs row operations that zero column $k$ below the diagonal. Gram--Schmidt does the opposite: it \emph{right}-multiplies by a sequence of upper triangular matrices in order to reach \emph{orthonormal columns}:
\begin{equation}
  AR_1R_2 \cdots R_n = \hat{Q}.
\end{equation}
Right multiplication performs column operations, and column operations are exactly what Gram--Schmidt does: normalize one column, then subtract its component from the other columns.

\subsection{The Step Matrices}

Write $\vec{a}_j^{(k-1)}$ for the $j$-th column after $k - 1$ steps, i.e., $\vec{a}_j$ with its components along $\vec{q}_1, \dots, \vec{q}_{k-1}$ removed, and $\vec{a}_j^{(0)} = \vec{a}_j$. At the start of step $k$ the matrix is
\begin{equation}
  \begin{bmatrix} \vec{q}_1 & \cdots & \vec{q}_{k-1} & \vec{a}_k^{(k-1)} & \cdots & \vec{a}_n^{(k-1)} \end{bmatrix}:
\end{equation}
the first $k - 1$ columns are orthonormal and the remaining columns are orthogonal to them, but not yet to each other. Let
\begin{equation}
  r_{kk} \doteq \norm{\vec{a}_k^{(k-1)}}, \qquad \vec{q}_k \doteq \vec{a}_k^{(k-1)}/r_{kk}, \qquad r_{kj} \doteq \ip{\vec{q}_k}{\vec{a}_j^{(k-1)}} \quad (j > k).
\end{equation}
The matrix $R_k$ equals the identity except in row $k$:
\begin{equation}
  R_k = \begin{bmatrix}
    1 & & & & & \\
    & \ddots & & & & \\
    & & \frac{1}{r_{kk}} & -\frac{r_{k,k+1}}{r_{kk}} & \cdots & -\frac{r_{kn}}{r_{kk}} \\
    & & & 1 & & \\
    & & & & \ddots & \\
    & & & & & 1
  \end{bmatrix},
\end{equation}
and right multiplication by it gives
\begin{equation}
  \begin{bmatrix} \vec{q}_1 & \cdots & \vec{q}_{k-1} & \vec{a}_k^{(k-1)} & \cdots & \vec{a}_n^{(k-1)} \end{bmatrix}R_k = \begin{bmatrix} \vec{q}_1 & \cdots & \vec{q}_k & \vec{a}_{k+1}^{(k)} & \cdots & \vec{a}_n^{(k)} \end{bmatrix},
\end{equation}
with $\vec{a}_j^{(k)} = \vec{a}_j^{(k-1)} - r_{kj}\vec{q}_k$. Indeed, column $k$ of the product is $\vec{a}_k^{(k-1)}/r_{kk} = \vec{q}_k$, and column $j > k$ is $\vec{a}_j^{(k-1)} - \vec{a}_k^{(k-1)}\,r_{kj}/r_{kk} = \vec{a}_j^{(k-1)} - r_{kj}\vec{q}_k$. Note the factor $1/r_{kk}$ in the off-diagonal entries: the coefficient multiplies the \emph{unnormalized} column $\vec{a}_k^{(k-1)}$, not $\vec{q}_k$. For $k = 1$,
\begin{equation}
  (R_1)_{1j} = -\frac{\ip{\vec{q}_1}{\vec{a}_j}}{\norm{\vec{a}_1}} = -\frac{\ip{\vec{a}_1}{\vec{a}_j}}{\norm{\vec{a}_1}^2}.
\end{equation}

\subsection{The Product Form Is Modified Gram--Schmidt}

Whether an algorithm is CGS or MGS is decided by the inner product defining the projection coefficients: CGS uses $\ip{\vec{q}_k}{\vec{a}_j}$ with the original column, MGS uses $\ip{\vec{q}_k}{\vec{a}_j^{(k-1)}}$ with the updated one. The product form uses $\vec{a}_j^{(k-1)}$, so it is \emph{MGS}. Structurally it cannot be anything else: $R_k$ acts on the \emph{current} matrix, whose $j$-th column is already $\vec{a}_j^{(k-1)}$, the original $\vec{a}_j$ having been overwritten. Each step normalizes column $k$ and immediately removes its component from all columns to the right, which is the row-oriented loop order of MGS.

\subsection{The Reduced QR Factorization}

Each $R_k$ is upper triangular, and so is their product:
\begin{equation}
  R_1R_2 \cdots R_n = \hat{R}^{-1} \quad\Longrightarrow\quad A = \hat{Q}\hat{R}.
\end{equation}
Here $\hat{Q} \in \K^{m \times n}$ is the factor of the \emph{reduced QR factorization}: it has only $n$ orthonormal columns, $\hat{Q}^\ast\hat{Q} = I_n$, but $\hat{Q}\hat{Q}^\ast \ne I_m$ (it is the orthogonal projector onto $\range(A)$). The factor $\hat{R} \in \K^{n \times n}$ is upper triangular and invertible.

In complete analogy with LU, $R_k^{-1}$ is also nontrivial only in row $k$, and that row is simply
\begin{equation}
  \begin{bmatrix} r_{kk} & r_{k,k+1} & \cdots & r_{kn} \end{bmatrix},
\end{equation}
with no reciprocals and no minus signs. Consequently
\begin{equation}
  \hat{R} = R_n^{-1} \cdots R_1^{-1} = \begin{bmatrix} r_{11} & r_{12} & \cdots & r_{1n} \\ & r_{22} & \cdots & r_{2n} \\ & & \ddots & \vdots \\ & & & r_{nn} \end{bmatrix}
\end{equation}
is obtained by stacking the $k$-th rows of the $R_k^{-1}$, just as $L = L_1^{-1} \cdots L_{m-1}^{-1}$ collects the multipliers column by column (\Cref{thm:elem-inverse}).

\section{Error Amplification in Products}

Consider algorithms of the form
\begin{equation}
  A_k = X_kA_{k-1}, \qquad A_0 = A.
\end{equation}
LU ($X_k = L_k$) and Householder QR ($X_k$ orthogonal, below) are of this type; Gram--Schmidt multiplies from the right and is analyzed symmetrically. Every step incurs rounding errors. How stable is such an algorithm?

\subsection{Forward Error of One Step}

One can show that floating-point matrix multiplication satisfies
\begin{equation}
  \fl(XA) = XA + E + O(\epsm^2), \qquad \norm{E}_2 \le O(\epsm)\norm{X}_2\norm{A}_2.
\end{equation}
The error $E$ is proportional to $\norm{X}\norm{A}$, \emph{not} to $\norm{XA}$. If $X$ and $A$ are both large but their product is small because of cancellation, the error is still measured against the large quantities and can be large relative to the result. This is the same phenomenon as the error of LU being governed by $\abs{L}\abs{U}$ rather than $\abs{A}$. The $O(\epsm)$ hides a dimension-dependent constant, and $O(\epsm^2)$ signals a first-order analysis.

\subsection{Rewriting as a Backward Error}

If $X$ is invertible, we can factor $X$ out of $E = X \cdot X^{-1}E$:
\begin{equation}
  \fl(XA) = X(A + F) + O(\epsm^2), \qquad F \doteq X^{-1}E.
\end{equation}
The computed result is the \emph{exact product with a perturbed input $A + F$}; the error has been pushed back onto the input, which is the backward-error viewpoint. The perturbation satisfies
\begin{equation}
  \norm{F}_2 \le \norm{X^{-1}}_2\norm{E}_2 \le O(\epsm)\,\kappa(X)\,\norm{A}_2,
\end{equation}
i.e., the relative backward error is $\norm{F}_2/\norm{A}_2 \le O(\epsm)\kappa(X)$.

\begin{theorem}[The Transformation Determines the Backward Error]
  \label{thm:product-error}
  For one step $A \mapsto XA$ of a product-form algorithm with invertible $X$, the computed result equals $X(A + F)$ with
  \begin{equation}
    \frac{\norm{F}_2}{\norm{A}_2} \le O(\epsm)\,\kappa(X).
  \end{equation}
\end{theorem}

Two condition numbers must be kept apart:
\begin{itemize}
  \item $\kappa(X)$ is the condition number of the transformation chosen by the \emph{algorithm}. It determines the backward error, i.e., whether the algorithm itself is stable, and the algorithm designer controls it.
  \item $\kappa(A)$ is the condition number of the \emph{problem}. It is fixed by the input and determines how much the backward error is amplified into forward error (forward error $\lesssim \kappa(A) \times$ backward error).
\end{itemize}
The design principle follows: choose transformations with small condition numbers. Orthogonal matrices have $\kappa = 1$, the best possible.

\section{Lower Triangularization: The Case of LU}

Apply the framework to LU. After $k$ steps of elimination the accumulated transformation is $X = L_k \cdots L_1$, and its inverse is the unit lower triangular matrix of multipliers,
\begin{equation}
  L = X^{-1} = \begin{bmatrix} 1 & & & \\ \ell_{21} & 1 & & \\ \vdots & \ddots & \ddots & \\ \ell_{m1} & \cdots & \ell_{m,m-1} & 1 \end{bmatrix}.
\end{equation}
Since $\kappa(X) = \kappa(X^{-1}) = \kappa(L)$, \Cref{thm:product-error} gives a backward error of about $O(\epsm)\kappa(L)\norm{A}_2$. The stability of LU thus comes down to the condition number of $L$. Partial pivoting chooses the largest entry of each column as pivot and so guarantees $\abs{\ell_{ij}} \le 1$.

\subsection{$\norm{L}$ Is Under Control}

All entries of $L$ have absolute value at most $1$, so its norm is only polynomial in $m$:
\begin{equation}
  \norm{L}_2 \le \norm{L}_F \le \sqrt{\frac{m(m+1)}{2}} \le m.
\end{equation}
Here $\norm{\cdot}_F$ is the Frobenius norm (see \Cref{sec:frobenius} in \Cref{app:norms}): $L$ has $m$ ones on the diagonal and $m(m-1)/2$ entries of absolute value at most $1$ below it, so $\norm{L}_F^2 \le m + m(m-1)/2 = m(m+1)/2$. The bound is loose, but all that matters is that it is polynomial rather than exponential in $m$.

\subsection{$\norm{L^{-1}}$ Is Not}

Controlling $\kappa(L) = \norm{L}_2\norm{L^{-1}}_2$ also requires $\norm{L^{-1}}$ to be moderate, and $\abs{\ell_{ij}} \le 1$ does \emph{not} imply this. Inverting $L$ (forward substitution) uses the results of all previous rows in every row, so values pile up and the entries of $L^{-1}$ can grow exponentially.

\begin{example}[Exponential Growth of $L^{-1}$]
  \label{ex:linv}
  The factor $L$ of \Cref{ex:growth} has $-1$ everywhere below the diagonal, so $\abs{\ell_{ij}} \le 1$, yet
  \begin{equation}
    L = \begin{bmatrix} 1 & & & \\ -1 & 1 & & \\ -1 & -1 & 1 & \\ -1 & -1 & -1 & 1 \end{bmatrix}, \qquad
    L^{-1} = \begin{bmatrix} 1 & & & \\ 1 & 1 & & \\ 2 & 1 & 1 & \\ 4 & 2 & 1 & 1 \end{bmatrix}.
  \end{equation}
  For general $m$, $(L^{-1})_{ij} = 2^{i-j-1}$ for $i > j$, and the bottom-left entry is $2^{m-2}$, so
  \begin{equation}
    \norm{L^{-1}}_2 \ge 2^{m-2}, \qquad \kappa(L) \ge 2^{m-2}.
  \end{equation}
  To see where the doubling comes from, note that computing $L^{-1}$ means solving $L\vec{x} = \vec{b}$. By forward substitution, the $i$-th equation $x_i - x_1 - \cdots - x_{i-1} = b_i$ gives
  \begin{equation}
    x_i = b_i + (x_1 + \cdots + x_{i-1}).
  \end{equation}
  For the first column, $\vec{b} = \vec{e}_1$: $x_1 = 1$, $x_2 = 1$, $x_3 = 1 + 1 = 2$, $x_4 = 1 + 1 + 2 = 4$, and so on. Each entry is the sum of all previous ones, so from the second entry on they double, $x_i = 2^{i-2}$. The other columns are the same pattern shifted down.
\end{example}

\subsection{Large $L^{-1}$ Causes Large $U$}

From $A = LU$,
\begin{equation}
  U = L^{-1}A,
\end{equation}
so a large $\norm{L^{-1}}$ can produce a large $U$: \emph{large $U$ comes from large $L^{-1}$}. This connects the condition-number viewpoint with the growth factor of \Cref{ch:pivoting,ch:ls}; they describe the same thing. In \Cref{ex:growth} the last column of $A$ is all ones, so the last column of $U$ is $L^{-1}$ times the all-ones vector, i.e., the row sums of $L^{-1}$:
\begin{equation}
  1 + \sum_{j=1}^{i-1}2^{i-j-1} = 1 + (2^{i-1} - 1) = 2^{i-1},
\end{equation}
which is exactly the last column $1, 2, 4, 8, \dots$ of $U$.

Partial pivoting thus solves only half of the problem: it controls $\norm{L}$ but not $\norm{L^{-1}}$, and in the worst case LU with partial pivoting can still be unstable. Such matrices are extremely rare, so in practice it is used as a stable algorithm; but to \emph{guarantee} well-conditioned transformations we must switch to orthogonal ones.

\section{Gram--Schmidt Is Also a Lower Triangularization}

The upper triangular $R_k$ is nontrivial only in row $k$, so its transpose $L_k^\T \doteq R_k$ has $L_k$ nontrivial only in column $k$: an elementary lower triangular matrix (with diagonal entry $1/r_{kk}$ instead of $1$). Then $AR_1 \cdots R_n = \hat{Q}$ reads $AL_1^\T \cdots L_n^\T = \hat{Q}$, and transposing,
\begin{equation}
  L_n \cdots L_1A^\T = \hat{Q}^\T.
\end{equation}
Gram--Schmidt on $A$ is therefore the same as left-multiplying $A^\T$ by a sequence of lower triangular matrices. It too falls into the error-amplification framework, and its stability is likewise governed by the condition numbers of these triangular factors.

\section{Householder Triangularization}

\subsection{Idea: Multiply by Orthogonal Matrices}

LU and Gram--Schmidt suffer from the same defect: each step applies a triangular matrix, whose condition number is not under control. Following the principle ``choose well-conditioned transformations'', the most direct remedy is to make every transformation orthogonal.

Householder triangularization, like LU, multiplies from the \emph{left} and aims for an \emph{upper triangular} matrix; the difference is that it applies unitary matrices $Q_k$. Each step zeros one column below the diagonal:
\begin{equation}
  \begin{bmatrix} \times & \times & \times \\ \times & \times & \times \\ \times & \times & \times \\ \times & \times & \times \\ \times & \times & \times \end{bmatrix}
  \xrightarrow{Q_1}
  \begin{bmatrix} \times & \times & \times \\ 0 & \times & \times \\ 0 & \times & \times \\ 0 & \times & \times \\ 0 & \times & \times \end{bmatrix}
  \xrightarrow{Q_2}
  \begin{bmatrix} \times & \times & \times \\ 0 & \times & \times \\ 0 & 0 & \times \\ 0 & 0 & \times \\ 0 & 0 & \times \end{bmatrix}
  \xrightarrow{Q_3}
  \begin{bmatrix} \times & \times & \times \\ 0 & \times & \times \\ 0 & 0 & \times \\ 0 & 0 & 0 \\ 0 & 0 & 0 \end{bmatrix}.
\end{equation}
After $n$ steps,
\begin{equation}
  Q_n \cdots Q_1A = R \quad\iff\quad A = \underbrace{(Q_n \cdots Q_1)^{-1}}_{Q}R = Q_1^\ast \cdots Q_n^\ast R.
\end{equation}
This is the \emph{full QR factorization}: $Q \in \K^{m \times m}$ is a square unitary matrix and $R \in \K^{m \times n}$ has $m - n$ zero rows at the bottom. Taking the first $n$ columns of $Q$ and the first $n$ rows of $R$ recovers the reduced QR factorization. In short, Gram--Schmidt uses triangular matrices to \emph{produce} an orthogonal matrix, whereas Householder uses orthogonal matrices to \emph{produce} a triangular one. It remains to choose the $Q_k$.

\subsection{The Householder Reflector}

$Q_1$ must map the first column $\vec{a}_1$ to a multiple of $\vec{e}_1$. Unitary maps preserve the $2$-norm, so there are only two possible targets,
\begin{equation}
  Q_1\vec{a}_1 = \pm\norm{\vec{a}_1}\vec{e}_1.
\end{equation}
Take the $+$ sign first. Since $\vec{a}_1$ and $\norm{\vec{a}_1}\vec{e}_1$ have the same length, they are mirror images across the perpendicular bisector of the segment joining them. This bisector is the hyperplane through the origin with normal vector
\begin{equation}
  \vec{v} = \vec{a}_1 - \norm{\vec{a}_1}\vec{e}_1,
\end{equation}
and \emph{reflecting} across it sends $\vec{a}_1$ to $\norm{\vec{a}_1}\vec{e}_1$.

\begin{center}
\begin{tikzpicture}[scale=1.4]
  \draw[-{Latex}] (-0.3,0) -- (3.4,0) node[right] {$\vec{e}_1$};
  \coordinate (O) at (0,0);
  \coordinate (X) at ({3*cos(50)},{3*sin(50)});
  \coordinate (E) at (3,0);
  \draw[dashed] (O) -- ({3.4*cos(25)},{3.4*sin(25)}) node[right] {$H$};
  \draw[-{Latex},thick] (O) -- (X) node[above] {$\vec{x}$};
  \draw[-{Latex},thick] (O) -- (E) node[below] {$\norm{\vec{x}}\vec{e}_1$};
  \draw[-{Latex},thick,red!70!black] (X) -- (E) node[midway,right] {$\norm{\vec{x}}\vec{e}_1 - \vec{x}$};
\end{tikzpicture}
\end{center}

\begin{definition}[Householder Reflector]
  For a nonzero $\vec{v} \in \K^m$, the \emph{Householder reflector} is
  \begin{equation}
    F = I - 2\frac{\vec{v}\vec{v}^\ast}{\vec{v}^\ast\vec{v}}.
  \end{equation}
\end{definition}

$F$ flips the component of a vector along $\vec{v}$ and leaves the component orthogonal to $\vec{v}$ unchanged:
\begin{itemize}
  \item $\vec{u} \perp \vec{v} \implies F\vec{u} = \vec{u}$: vectors in the hyperplane stay fixed;
  \item $\vec{u} \parallel \vec{v} \implies F\vec{u} = -\vec{u}$: the normal direction is flipped.
\end{itemize}
Clearly $F$ is Hermitian, and reflecting twice returns every vector to where it started, $F^2 = I$. Hence $F^\ast F = F^2 = I$: $F$ is unitary and is its own inverse.

We verify $Q_1\vec{a}_1 = \norm{\vec{a}_1}\vec{e}_1$ in the real case, with $Q_1 = I - 2\vec{v}\vec{v}^\ast/(\vec{v}^\ast\vec{v})$:
\begin{equation}
  \vec{v}^\ast\vec{v} = 2\norm{\vec{a}_1}^2 - 2\norm{\vec{a}_1}a_{11}, \qquad \vec{v}^\ast\vec{a}_1 = \norm{\vec{a}_1}^2 - \norm{\vec{a}_1}a_{11} = \tfrac{1}{2}\vec{v}^\ast\vec{v},
\end{equation}
so
\begin{equation}
  Q_1\vec{a}_1 = \vec{a}_1 - 2\vec{v}\frac{\vec{v}^\ast\vec{a}_1}{\vec{v}^\ast\vec{v}} = \vec{a}_1 - \vec{v} = \norm{\vec{a}_1}\vec{e}_1.
\end{equation}

\subsection{Choosing the Sign}

The targets $+\norm{\vec{a}_1}\vec{e}_1$ and $-\norm{\vec{a}_1}\vec{e}_1$ correspond to two different (mutually perpendicular) reflection hyperplanes. Mathematically either one works, but in finite precision one should choose the target \emph{farther from $\vec{a}_1$}:
\begin{equation}
  Q_1\vec{a}_1 = -\sign(a_{11})\norm{\vec{a}_1}\vec{e}_1 \quad\implies\quad \vec{v} \propto \sign(a_{11})\norm{\vec{a}_1}\vec{e}_1 + \vec{a}_1.
\end{equation}
The reason is cancellation. Suppose $\vec{a}_1$ is nearly parallel to $\vec{e}_1$ with $a_{11} > 0$. With the $+$ sign,
\begin{equation}
  \norm{\vec{v}_+} = \norm{\norm{\vec{a}_1}\vec{e}_1 - \vec{a}_1} \ll \norm{\vec{a}_1},
\end{equation}
so $\vec{v}_+$ is the difference of two nearly equal vectors and its first entry $\norm{\vec{a}_1} - a_{11}$ suffers catastrophic cancellation. The computed $\norm{\vec{a}_1}$ carries a relative error $O(\epsm)$, i.e., an absolute error of about $\epsm\norm{\vec{a}_1}$, which is large relative to the small $\norm{\vec{v}_+}$; the direction of $\vec{v}$, and with it the reflection, becomes inaccurate. With $-\sign(a_{11})$, the first entry of $\vec{v}$ is $a_{11} + \sign(a_{11})\norm{\vec{a}_1}$, a sum of two numbers of the same sign with no cancellation, and $\abs{v_1} \ge \norm{\vec{a}_1}$, so $\norm{\vec{v}}$ is never smaller than $\norm{\vec{a}_1}$. We take $\sign(0) = 1$.


\begin{center}
\begin{minipage}[t]{0.36\textwidth}
\centering
\begin{tikzpicture}[scale=1.2]
  \def\R{2.2}\def\th{24}\def\r{0.3}
  \coordinate (O) at (0,0);
  \coordinate (A) at ({\R*cos(\th)},{\R*sin(\th)});
  \coordinate (T) at (\R,0);
  \draw[-{Latex}] (-0.3,0) -- (3.1,0) node[right] {$\vec{e}_1$};
  \draw[dashed] (O) -- ({3.0*cos(\th/2)},{3.0*sin(\th/2)}) node[right] {$H_+$};
  \fill[red!15] (T) circle (\r);
  \draw[-{Latex},thick] (O) -- (A) node[above left] {$\vec{a}_1$};
  \draw[-{Latex},thick] (O) -- (T);
  \node[below=4pt] at (T) {$\norm{\vec{a}_1}\vec{e}_1$};
  \draw[-{Latex},thin,dashed,red!60!black] ($(T)+(0,\r)$) -- (A);
  \draw[-{Latex},thin,dashed,red!60!black] ($(T)+(\r,0)$) -- (A);
  \draw[-{Latex},thick,red!70!black] (T) -- (A);
  \node[red!70!black] at (2.45,0.8) {$\vec{v}_+$};
\end{tikzpicture}

\small Reflecting onto $+\norm{\vec{a}_1}\vec{e}_1$: $\vec{v}_+$ is short, so its direction is unreliable.
\end{minipage}\hfill
\begin{minipage}[t]{0.58\textwidth}
\centering
\begin{tikzpicture}[scale=1.2]
  \def\R{2.2}\def\th{24}\def\r{0.3}
  \coordinate (O) at (0,0);
  \coordinate (A) at ({\R*cos(\th)},{\R*sin(\th)});
  \coordinate (T) at (-\R,0);
  \draw[-{Latex}] (-2.8,0) -- (3.1,0) node[right] {$\vec{e}_1$};
  \draw[dashed] ({-0.8*cos(90+\th/2)},{-0.8*sin(90+\th/2)}) -- ({1.3*cos(90+\th/2)},{1.3*sin(90+\th/2)}) node[above] {$H_-$};
  \fill[red!15] (T) circle (\r);
  \draw[-{Latex},thick] (O) -- (A) node[above right] {$\vec{a}_1$};
  \draw[-{Latex},thick] (O) -- (T);
  \node[below=4pt] at (T) {$-\norm{\vec{a}_1}\vec{e}_1$};
  \draw[-{Latex},thin,dashed,red!60!black] ($(T)+(0,\r)$) -- (A);
  \draw[-{Latex},thin,dashed,red!60!black] ($(T)+(\r,0)$) -- (A);
  \draw[-{Latex},thick,red!70!black] (T) -- (A);
  \node[red!70!black,above=3pt] at ($(T)!0.3!(A)$) {$\vec{v}$};
\end{tikzpicture}

\small Reflecting onto $-\norm{\vec{a}_1}\vec{e}_1$: $\vec{v}$ is long, so its direction is reliable.
\end{minipage}
\end{center}

\subsection{Step $k$: The Block Form}

For $k > 1$, the first $k - 1$ rows and columns are finished and must not change. So we take
\begin{equation}
  Q_k = \begin{bmatrix} I_{k-1} & 0 \\ 0 & F_k \end{bmatrix},
\end{equation}
where $F_k$ is the $(m - k + 1) \times (m - k + 1)$ Householder reflector built as in step $1$ from the first column $\vec{x} = A[k{:}m, k]$ of the trailing submatrix $A[k{:}m, k{:}n]$. Earlier results are preserved: the block $I_{k-1}$ leaves the first $k - 1$ rows untouched, and the first $k - 1$ columns are already zero from row $k$ down, and $F_k$ maps zero vectors to zero. The matrix $Q_k$ is itself a Householder reflector, with normal vector $\vec{v}$ padded by $k - 1$ leading zeros, and is unitary.

\begin{algorithm}[H]
  \caption{Householder QR, in place.}
  \label{alg:householder}
  \begin{algorithmic}[1]
    \For{$k = 1, \dots, n$}
      \State $\vec{x} \gets A[k{:}m, k]$
      \State $\vec{v} \gets \sign(x_1)\norm{\vec{x}}_2\vec{e}_1 + \vec{x}$
      \State $\vec{v} \gets \vec{v}/v_1$ \Comment{scale so that $v_1 = 1$}
      \State $A[k{:}m, k{:}n] \gets A[k{:}m, k{:}n] - 2\vec{v}\left(\vec{v}^\ast A[k{:}m, k{:}n]\right)/(\vec{v}^\ast\vec{v})$
      \State $A[k{+}1{:}m, k] \gets \vec{v}[2{:}\mathrm{end}]$ \Comment{store $\vec{v}$ below the diagonal}
    \EndFor
    \State \Return $A$ \Comment{upper triangle: $R$; strictly lower triangle: the vectors $\vec{v}$}
  \end{algorithmic}
\end{algorithm}

Some remarks:
\begin{itemize}
  \item \emph{$F_k$ is never formed.} Following the parentheses, first compute the row vector $\vec{v}^\ast A[k{:}m, k{:}n]$, then perform a rank-one update. Each step costs $O((m - k)(n - k))$; forming $\vec{v}\vec{v}^\ast$ first would cost an order of magnitude more.
  \item \emph{Scaling $\vec{v}$ by $1/v_1$} makes its first entry $1$. The reflector depends only on the direction of $\vec{v}$, so scaling does not change $F_k$. Since $\abs{v_1} \ge \norm{\vec{x}} > 0$ (for $\vec{x} \ne \zerovec$), this never divides by zero.
  \item \emph{Storage.} After the update, $A[k{+}1{:}m, k]$ should be zero, which frees exactly enough room for the remaining $m - k$ entries of $\vec{v}$; the leading $v_1 = 1$ is implicit.
  \item On exit the upper triangle of $A$ is $R$, with diagonal entries $r_{kk} = -\sign(x_1)\norm{\vec{x}}$. The matrix $Q$ is never formed explicitly; it is stored implicitly through the vectors $\vec{v}$ (see \Cref{ch:qr}).
  \item $\sign(0)$ must be taken as $1$; otherwise $x_1 = 0$ gives $\vec{v} = \vec{x}$, which is wrong.
\end{itemize}

\subsection{Back into the Framework}

Each $X_k = Q_k$ is unitary, so $\norm{Q_k}_2 = \norm{Q_k^{-1}}_2 = 1$ and $\kappa(Q_k) = 1$. By \Cref{thm:product-error},
\begin{equation}
  \norm{F}_2 \le O(\epsm)\,\kappa(Q_k)\,\norm{A}_2 = O(\epsm)\norm{A}_2:
\end{equation}
the transformation does not amplify the error at all (see \Cref{app:norms} for the properties of unitary matrices used here). Moreover, unitary matrices preserve the $2$-norm, so the intermediate matrices $A_k$ all have norm $\norm{A}_2$, and there is no element growth as in LU. Accumulating over all steps (see \Cref{sec:hh-supp}), the total backward error grows only linearly with the number of steps: Householder QR is backward stable. This presumes that the $Q_k$ are computed accurately, i.e., that the direction of $\vec{v}$ is accurate, which is exactly what the sign choice ensures.

\section{(SUPPLEMENT) Further Topics}
\label{sec:hh-supp}

\paragraph{Three factorizations as products.}
\begin{center}
\begin{tabular}{lllll}
  \toprule
  & Side & Transformation & Target & Condition of transformation \\
  \midrule
  LU & left & lower triangular $L_k$ & upper triangular $U$ & can be large \\
  Gram--Schmidt & right & upper triangular $R_k$ & orthonormal $\hat{Q}$ & depends on $A$ \\
  Householder & left & unitary $Q_k$ & upper triangular $R$ & $1$ \\
  \bottomrule
\end{tabular}
\end{center}
Trefethen and Bau call Gram--Schmidt \emph{triangular orthogonalization} and Householder \emph{orthogonal triangularization}~\cite[Lectures 8 and 10]{TB}.

\paragraph{The condition of the Gram--Schmidt transformation.} The accumulated transformation of Gram--Schmidt is $R_1 \cdots R_n = \hat{R}^{-1}$. Since $A = \hat{Q}\hat{R}$ with $\hat{Q}$ having orthonormal columns, $A$ and $\hat{R}$ have the same singular values, so $\kappa(\hat{R}) = \kappa(A)$. The triangular transformations of Gram--Schmidt are thus exactly as ill-conditioned as the problem: the more ill-conditioned $A$, the more the error is amplified. This matches the loss of orthogonality $\approx \epsm\kappa(A)$ of MGS from \Cref{ch:ls}.

\paragraph{Backward error over many steps.} Writing each step as $\tilde{A}_k = X_k(\tilde{A}_{k-1} + F_k)$ and pulling every perturbation back to the original input,
\begin{equation}
  \tilde{A}_n = X_n \cdots X_1\left(A + \sum_{k=1}^{n}(X_{k-1} \cdots X_1)^{-1}F_k\right) + O(\epsm^2).
\end{equation}
When every $X_k$ is unitary, so is $(X_{k-1} \cdots X_1)^{-1}$, which does not change the norm of $F_k$; moreover $\norm{\tilde{A}_{k-1}}_2 = \norm{A}_2$ to first order, so each $\norm{F_k}_2 \le O(\epsm)\norm{A}_2$ and
\begin{equation}
  \norm{F_{\mathrm{total}}}_2 \lesssim O(n\epsm)\norm{A}_2.
\end{equation}
This is the argument behind Demmel's Lemma 3.1 and Theorem 3.5~\cite[\S3.4.3]{Demmel}, and Trefethen and Bau state the result as $\tilde{Q}\tilde{R} = A + \delta A$ with $\norm{\delta A}/\norm{A} = O(\epsm)$~\cite[Theorem 16.1]{TB}, where $\tilde{Q}$ is the \emph{exactly} unitary product of the reflectors defined by the computed vectors $\tilde{\vec{v}}_k$.

\paragraph{Orthogonality of $Q$.} The $\hat{Q}$ computed by MGS loses orthogonality by about $O(\epsm\kappa(A))$, whereas the $Q$ given implicitly by Householder is orthogonal up to $O(\epsm)$, independently of $\kappa(A)$, consistent with the table above.

\paragraph{Operation counts.} MGS costs about $2mn^2$ flops. Householder QR costs about $2mn^2 - \frac{2}{3}n^3$ flops: step $k$ updates an $(m - k + 1) \times (n - k + 1)$ block with about $4$ flops per entry, and
\begin{equation}
  \sum_{k=1}^{n}4(m - k + 1)(n - k + 1) \approx 4\int_0^n (m - t)(n - t)\,\mathrm{d}t = 2mn^2 - \frac{2}{3}n^3.
\end{equation}
For $m = n$ Householder costs $\frac{4}{3}n^3$, twice as much as LU ($\frac{2}{3}n^3$); this is the price of stability.

\paragraph{Using $Q$ implicitly for least squares.} To solve $\min_{\vec{x}}\norm{A\vec{x} - \vec{b}}_2$, we never form $Q$; we apply the stored $\vec{v}_k$ to $\vec{b}$ in turn to obtain $Q^\ast\vec{b}$, then solve $\hat{R}\vec{x} = (Q^\ast\vec{b})[1{:}n]$ by back substitution. The norm of the remaining part $(Q^\ast\vec{b})[n{+}1{:}m]$ is the minimal residual. Details are in \Cref{ch:qr}.

\paragraph{The complex case.} For complex data $\vec{v}^\ast\vec{a}_1$ is in general not real, so the verification above does not carry over verbatim. One takes $\sign(x_1) \doteq x_1/\abs{x_1}$ (and $1$ if $x_1 = 0$) and $\vec{v} = \sign(x_1)\norm{\vec{x}}\vec{e}_1 + \vec{x}$. Then the first entry of $\vec{v}$ adds two numbers with the same phase, again avoiding cancellation; one checks $\vec{v}^\ast\vec{x} = \norm{\vec{x}}\abs{x_1} + \norm{\vec{x}}^2 = \frac{1}{2}\vec{v}^\ast\vec{v}$, so $F\vec{x} = \vec{x} - \vec{v} = -\sign(x_1)\norm{\vec{x}}\vec{e}_1$.

\chapter{Using $Q$ Implicitly, Givens Rotations, and Existence and Uniqueness of QR}
\label{ch:qr}

Supplementary reading:
\begin{itemize}
  \item \cite{TB} Lectures 7 and 10.
  \item \cite{Demmel} Sections 3.2.2 and 3.4.
\end{itemize}

\section{Picking Up from Householder QR}

At the end of Householder QR (\Cref{alg:householder}), the upper triangle of $A$ holds $R$ and the strictly lower triangle holds the Householder vectors $\vec{v}_k$; $Q$ is never computed explicitly. This raises two questions:
\begin{itemize}
  \item Applications need $Q^\ast\vec{b}$ (for least squares) or $Q\vec{x}$. How do we compute them from the $\vec{v}_k$ alone?
  \item Every Householder step processes a whole column. If $A$ is already close to upper triangular and only a few entries need to be eliminated, is this wasteful?
\end{itemize}
The first leads to the implicit use of $Q$, the second to Givens rotations. Finally we return to theory: does a QR factorization always exist, and is it unique?

\section{Using $Q$}

\subsection{Implicit Storage}

Householder QR does not store $Q$ as a matrix, only the vectors $\vec{v}_k$ defining the reflectors
\begin{equation}
  Q_k = I - 2\frac{\vec{v}_k\vec{v}_k^\ast}{\vec{v}_k^\ast\vec{v}_k}.
\end{equation}
Here $\vec{v}_k$ is regarded as a vector of length $m$ whose first $k - 1$ entries are zero, so $Q_k$ acts only on rows $k$ through $m$; this is the same as the block form $Q_k = \diag(I_{k-1}, F_k)$ of \Cref{ch:householder}. From $Q_n \cdots Q_1A = R$, and since each $Q_k$ is Hermitian and its own inverse ($Q_k^\ast = Q_k = Q_k^{-1}$),
\begin{equation}
  Q = (Q_n \cdots Q_1)^{-1} = Q_1Q_2 \cdots Q_n, \qquad Q^\ast = Q_n \cdots Q_1.
\end{equation}

\subsection{Two Algorithms}

In $Q^\ast\vec{b} = Q_n \cdots Q_1\vec{b}$ the first factor to act is $Q_1$, so $k$ runs upward from $1$; in $Q\vec{x} = Q_1 \cdots Q_n\vec{x}$ the first factor to act is $Q_n$, so $k$ runs downward from $n$.

\begin{algorithm}[H]
  \caption{Implicit computation of $Q^\ast\vec{b}$.}
  \label{alg:qtb}
  \begin{algorithmic}[1]
    \For{$k = 1, \dots, n$}
      \State $\vec{b}[k{:}m] \gets \vec{b}[k{:}m] - 2\vec{v}_k\left(\vec{v}_k^\ast\vec{b}[k{:}m]\right)/(\vec{v}_k^\ast\vec{v}_k)$
    \EndFor
  \end{algorithmic}
\end{algorithm}

\begin{algorithm}[H]
  \caption{Implicit computation of $Q\vec{x}$.}
  \label{alg:qx}
  \begin{algorithmic}[1]
    \For{$k = n, n - 1, \dots, 1$}
      \State $\vec{x}[k{:}m] \gets \vec{x}[k{:}m] - 2\vec{v}_k\left(\vec{v}_k^\ast\vec{x}[k{:}m]\right)/(\vec{v}_k^\ast\vec{v}_k)$
    \EndFor
  \end{algorithmic}
\end{algorithm}

As in \Cref{ch:householder}, following the parentheses we first compute the scalar $\vec{v}_k^\ast\vec{b}[k{:}m]$ and then update a vector, so each step costs only $O(m - k)$ and the total is about $4mn - 2n^2$ flops, the same order as a matrix--vector product. One should never form $\vec{v}_k\vec{v}_k^\ast$. The two algorithms differ only in the direction of the loop, much like forward and back substitution: the order of the factors in the product determines the order in which they act.

\subsection{Consequences of Implicit Storage}

\paragraph{In place.} Scaling $\vec{v}_k$ so that its $k$-th entry is $1$ makes that entry implicit, and the remaining $m - k$ entries fit into the positions below the diagonal of column $k$ that were just zeroed. The whole factorization needs no extra matrix storage and simply overwrites $A$.

\paragraph{A full QR factorization from only $n$ vectors.} $Q$ is $m \times m$, so storing it explicitly takes $m^2$ numbers; the implicit form takes only the roughly $mn - \frac{n^2}{2}$ positions below the diagonal of $A$. When $m \gg n$ the difference is large. If the reduced factor $\hat{Q}$ (the first $n$ columns of $Q$) is really needed, compute $Q\vec{e}_1, \dots, Q\vec{e}_n$ with \Cref{alg:qx}.

\paragraph{$Q$ is orthogonal by definition.} This point was stressed in the lecture. However much rounding error the computed $\vec{v}_k$ carry, $I - 2\vec{v}_k\vec{v}_k^\ast/(\vec{v}_k^\ast\vec{v}_k)$ is an exact reflector for \emph{every} nonzero $\vec{v}_k$, hence exactly unitary. Rounding errors only affect \emph{which} orthogonal matrix $Q$ is, not \emph{whether} it is orthogonal; they are absorbed into a backward perturbation of $A$. Contrast this with Gram--Schmidt, which computes every column of $\hat{Q}$ explicitly, so rounding errors directly destroy the orthogonality between columns, by about $\epsm\kappa(A)$ for MGS and $\epsm\kappa(A)^2$ for CGS. When the Householder $Q$ is used in a computation (for example applied to $\vec{b}$ as above), the only error comes from applying exact reflectors in floating point, and this step is backward stable.

\section{Givens Rotations}

\subsection{Motivation: Householder Can Be Overkill}

A Householder reflector processes a whole column at once: step $k$ performs a rank-one update of the whole block $A[k{:}m, k{:}n]$, no matter how many nonzeros lie below the diagonal in column $k$. If the matrix is already \emph{close to upper triangular}, with only a narrow band of nonzeros below the diagonal (for example an upper Hessenberg matrix, which has a single nonzero subdiagonal), each column needs only one or two entries eliminated and Householder is wasteful. We want an orthogonal transformation that zeros \emph{one entry at a time}.

\subsection{The $2 \times 2$ Case}

Consider the leading $2 \times 2$ block
\begin{equation}
  \begin{bmatrix} A_{11} & A_{12} \\ A_{21} & A_{22} \end{bmatrix}.
\end{equation}
We want to eliminate $A_{21}$. Geometrically, the first column $\vec{a}_1 = (A_{11}, A_{21})^\T$ is a vector in the plane, and it suffices to \emph{rotate} the plane so that $\vec{a}_1$ lands on the horizontal axis; $\vec{a}_2$ rotates along with it.

\begin{definition}[Givens Rotation]
  With $\norm{\vec{a}_1} = \sqrt{A_{11}^2 + A_{21}^2}$, the \emph{Givens rotation} for $\vec{a}_1$ is
  \begin{equation}
    G = \begin{bmatrix} c & s \\ -s & c \end{bmatrix}, \qquad c = \frac{A_{11}}{\norm{\vec{a}_1}}, \qquad s = \frac{A_{21}}{\norm{\vec{a}_1}}.
  \end{equation}
\end{definition}

A direct check gives
\begin{equation}
  G\begin{bmatrix} A_{11} \\ A_{21} \end{bmatrix} = \frac{1}{\norm{\vec{a}_1}}\begin{bmatrix} A_{11}^2 + A_{21}^2 \\ -A_{21}A_{11} + A_{11}A_{21} \end{bmatrix} = \begin{bmatrix} \norm{\vec{a}_1} \\ 0 \end{bmatrix}.
\end{equation}
Since $c^2 + s^2 = 1$, we have $G^\T G = I$ and $\det G = 1$, so $G$ is a rotation: if $\vec{a}_1$ makes the angle $\theta$ with the horizontal axis, then $c = \cos\theta$, $s = \sin\theta$, and $G$ rotates clockwise by $\theta$. Compare with Householder: a Householder reflector is a \emph{reflection} ($\det = -1$) that maps a whole vector onto the direction of $\vec{e}_1$ at once, whereas a Givens rotation is a \emph{rotation} ($\det = 1$) acting in a single two-dimensional plane and zeroing one entry.

\subsection{QR by Givens Rotations}

In the $m \times m$ setting, the Givens rotation acting on rows $i$ and $j$ is the identity except in four positions:
\begin{equation}
  G(i, j) = \begin{bmatrix} 1 & & & & \\ & c & & s & \\ & & \ddots & & \\ & -s & & c & \\ & & & & 1 \end{bmatrix},
\end{equation}
with $c$ in positions $(i, i)$ and $(j, j)$, $s$ in $(i, j)$, and $-s$ in $(j, i)$. Left multiplication by $G(i, j)$ changes only rows $i$ and $j$. To eliminate $A_{ji}$ ($j > i$), take
\begin{equation}
  c = \frac{A_{ii}}{\sqrt{A_{ii}^2 + A_{ji}^2}}, \qquad s = \frac{A_{ji}}{\sqrt{A_{ii}^2 + A_{ji}^2}},
\end{equation}
which is the $2 \times 2$ case transplanted to rows $i$ and $j$.

Existing zeros are preserved: when column $i$ is processed, the first $i - 1$ columns are already zero in both rows $i$ and $j$, and a linear combination of two zeros is zero. In column $i$ itself, row $i$ becomes $\sqrt{A_{ii}^2 + A_{ji}^2}$, row $j$ becomes $0$, and no other row changes. Eliminating the subdiagonal nonzeros one at a time until the matrix is upper triangular gives
\begin{equation}
  G_N \cdots G_1A = R, \qquad Q = G_1^\ast \cdots G_N^\ast.
\end{equation}
As with Householder, $Q$ is not formed explicitly; one stores the pairs $(c, s)$ of the rotations and applies them in sequence when needed, just as in \Cref{alg:qtb,alg:qx}. Every $G$ is orthogonal with $\kappa = 1$, so by the error-amplification framework of \Cref{ch:householder}, Givens QR is backward stable as well.

\subsection{Efficiency}

\paragraph{Dense matrices: slower than Householder.} Each rotation acts on two rows, at $6$ flops per column ($4$ multiplications and $2$ additions). Eliminating column $i$ takes $m - i$ rotations, each acting on rows of length about $n - i$, for a total of about
\begin{equation}
  \sum_{i=1}^{n}6(m - i)(n - i) \approx 3mn^2 - n^3 \text{ flops},
\end{equation}
about $1.5$ times the cost of Householder ($2mn^2 - \frac{2}{3}n^3$).

\paragraph{Sparse or nearly triangular matrices: possibly faster.} For an $n \times n$ upper Hessenberg matrix, only $n - 1$ rotations are needed (one per subdiagonal entry), each costing $O(n)$, for a total of about $3n^2$ flops. Givens touches only the entries that must be eliminated and so exploits sparsity naturally.

\section{Full and Reduced QR}

Let $A \in \K^{m \times n}$ with $m \ge n$.

\begin{definition}[Reduced and Full QR Factorizations]
  A \emph{reduced QR factorization} of $A$ is
  \begin{equation}
    A = \hat{Q}\hat{R}, \qquad \hat{Q} \in \K^{m \times n} \text{ with orthonormal columns}, \quad \hat{R} \in \K^{n \times n} \text{ upper triangular}.
  \end{equation}
  A \emph{full QR factorization} of $A$ is
  \begin{equation}
    A = QR, \qquad Q \in \K^{m \times m} \text{ unitary}, \quad R = \begin{bmatrix} \hat{R} \\ 0 \end{bmatrix} \in \K^{m \times n}, \ \hat{R} \text{ upper triangular}.
  \end{equation}
\end{definition}

Gram--Schmidt produces a reduced factorization; Householder and Givens produce a full one. Splitting $Q = \begin{bmatrix} \hat{Q} & \hat{Q}_\perp \end{bmatrix}$ by columns,
\begin{equation}
  QR = \begin{bmatrix} \hat{Q} & \hat{Q}_\perp \end{bmatrix}\begin{bmatrix} \hat{R} \\ 0 \end{bmatrix} = \hat{Q}\hat{R}.
\end{equation}
The extra $m - n$ columns $\hat{Q}_\perp$ multiply the zero block of $R$ and do not contribute to the product. Their role is to complete $\hat{Q}$ to an orthonormal basis of $\K^m$; when $A$ has full rank, the columns of $\hat{Q}_\perp$ span $\range(A)^\perp$.

\section{Existence}

\begin{theorem}[Existence of QR]
  \label{thm:qr-exists}
  Every $A \in \K^{m \times n}$ with $m \ge n$ has a full QR factorization and a reduced QR factorization.
\end{theorem}

The proof is constructive: run Gram--Schmidt, see where it can fail, and patch those cases.

\begin{proof}
  \emph{$A$ of full rank, reduced QR.} The only place where Gram--Schmidt can fail is the normalization $\vec{q}_k = \hat{\vec{q}}_k/\norm{\hat{\vec{q}}_k}$, which divides by zero if $\hat{\vec{q}}_k = \zerovec$. But
  \begin{equation}
    \norm{\hat{\vec{q}}_k} = 0 \iff \vec{a}_k \in \spn(\vec{q}_1, \dots, \vec{q}_{k-1}) = \spn(\vec{a}_1, \dots, \vec{a}_{k-1}),
  \end{equation}
  which contradicts full column rank. So Gram--Schmidt runs to completion, and all $r_{kk} = \norm{\hat{\vec{q}}_k} > 0$.

  \emph{$A$ rank deficient, reduced QR.} Run Gram--Schmidt as usual. Whenever $\hat{\vec{q}}_k = \zerovec$, set $r_{kk} = 0$ and take for $\vec{q}_k$ any unit vector orthogonal to $\vec{q}_1, \dots, \vec{q}_{k-1}$; such a vector exists because $k - 1 < n \le m$. Then still
  \begin{equation}
    \vec{a}_k = \sum_{i<k}r_{ik}\vec{q}_i + 0 \cdot \vec{q}_k,
  \end{equation}
  so $A = \hat{Q}\hat{R}$ holds, only with zeros on the diagonal of $\hat{R}$.

  \emph{Full QR.} Compute a reduced QR factorization, extend $\vec{q}_1, \dots, \vec{q}_n$ to an orthonormal basis of $\K^m$ by adding $m - n$ orthonormal vectors, and append $m - n$ zero rows to $\hat{R}$.
\end{proof}

\section{Uniqueness}

\subsection{When QR Is Not Unique}

\paragraph{Full QR with $m > n$ is not unique.} $\hat{Q}_\perp$ can be any orthonormal basis of $\range(\hat{Q})^\perp$, because it multiplies a zero block. For instance, replacing $\hat{Q}_\perp$ by $\hat{Q}_\perp U$ with any unitary $U \in \K^{(m-n) \times (m-n)}$ leaves the product unchanged.

\paragraph{Rank-deficient $A$.} When $\hat{\vec{q}}_k = \zerovec$, the choice of $\vec{q}_k$ is arbitrary.

\begin{example}[Non-Uniqueness for a Rank-Deficient Matrix]
  Let
  \begin{equation}
    A = \begin{bmatrix} 1 & 1 \\ 0 & 0 \\ 0 & 0 \end{bmatrix}, \qquad \vec{q}_1 = \vec{e}_1, \qquad \hat{\vec{q}}_2 = \vec{a}_2 - \vec{e}_1 = \zerovec.
  \end{equation}
  Both $\vec{q}_2 = \vec{e}_2$ and $\vec{q}_2 = \vec{e}_3$ work, with the same $\hat{R} = \begin{bmatrix} 1 & 1 \\ 0 & 0 \end{bmatrix}$.
\end{example}

\paragraph{Phases.} Even for full-rank $A$, given $\abs{d_k} = 1$ we may multiply $\vec{q}_k$ by $d_k$ and row $k$ of $\hat{R}$ by $\bar{d}_k$ without changing the product; in the real case, each column may flip its sign. Uniqueness therefore requires a normalization.

\subsection{Uniqueness of the Reduced QR Factorization}

\begin{theorem}[Uniqueness of Reduced QR]
  \label{thm:qr-unique}
  Let $A \in \K^{m \times n}$, $m \ge n$, have full column rank. Then the reduced QR factorization $A = \hat{Q}\hat{R}$ with $r_{jj} > 0$ for all $j$ is unique.
\end{theorem}

\begin{proof}
  Let $A = \hat{Q}\hat{R} = \tilde{Q}\tilde{R}$ be two such factorizations. We show $\tilde{\vec{q}}_k = \vec{q}_k$ column by column. Fix $k$, write $\vec{q} = \vec{q}_k$ and $\tilde{\vec{q}} = \tilde{\vec{q}}_k$, and use the inner product $\ip{\vec{x}}{\vec{y}} = \vec{x}^\ast\vec{y}$.

  First, what conditions do $\vec{q}$ and $\tilde{\vec{q}}$ satisfy? Full rank gives $r_{jj} \ne 0$, so $\hat{R}$ is upper triangular and invertible, and the first $k$ columns of $\hat{Q} = A\hat{R}^{-1}$ span the same space as the first $k$ columns of $A$:
  \begin{equation}
    \spn(\vec{q}_1, \dots, \vec{q}_k) = \spn(\vec{a}_1, \dots, \vec{a}_k) \qquad \text{for every } k.
  \end{equation}
  Hence $\vec{q} \in \spn(\vec{a}_1, \dots, \vec{a}_k)$ and $\vec{q}$ is orthogonal to $\vec{q}_1, \dots, \vec{q}_{k-1}$, hence to $\vec{a}_1, \dots, \vec{a}_{k-1}$. The same holds for $\tilde{\vec{q}}$.

  Since $\spn(\vec{a}_1, \dots, \vec{a}_k) = \spn(\vec{a}_1, \dots, \vec{a}_{k-1}, \vec{q})$, we can write
  \begin{equation}
    \tilde{\vec{q}} = \alpha\vec{q} + \sum_{j=1}^{k-1}\beta_j\vec{a}_j.
  \end{equation}
  Taking inner products with $\vec{q}$ and with $\tilde{\vec{q}}$, and using $\vec{q} \perp \vec{a}_j$ and $\tilde{\vec{q}} \perp \vec{a}_j$ for $j < k$,
  \begin{equation}
    \ip{\vec{q}}{\tilde{\vec{q}}} = \alpha, \qquad 1 = \ip{\tilde{\vec{q}}}{\tilde{\vec{q}}} = \alpha\ip{\tilde{\vec{q}}}{\vec{q}} = \alpha\bar{\alpha} = \abs{\alpha}^2.
  \end{equation}
  Therefore
  \begin{equation}
    \norm{\tilde{\vec{q}} - \alpha\vec{q}}^2 = \ip{\tilde{\vec{q}}}{\tilde{\vec{q}}} - \alpha\ip{\tilde{\vec{q}}}{\vec{q}} - \bar{\alpha}\ip{\vec{q}}{\tilde{\vec{q}}} + \abs{\alpha}^2 = 1 - 1 - 1 + 1 = 0,
  \end{equation}
  so $\tilde{\vec{q}} = \alpha\vec{q}$ with $\abs{\alpha} = 1$.

  Finally, $r_{kk} > 0$ fixes $\alpha$. From $A = \hat{Q}\hat{R}$ we get $r_{kk} = \ip{\vec{q}_k}{\vec{a}_k}$, and likewise $\tilde{r}_{kk} = \ip{\tilde{\vec{q}}_k}{\vec{a}_k} = \ip{\alpha\vec{q}_k}{\vec{a}_k} = \bar{\alpha}r_{kk}$. Both are positive real numbers, so $\bar{\alpha} = \tilde{r}_{kk}/r_{kk} > 0$, and with $\abs{\alpha} = 1$ this forces $\alpha = 1$.

  Thus $\tilde{\vec{q}}_k = \vec{q}_k$ for every $k$, i.e., $\tilde{Q} = \hat{Q}$, and then $\tilde{R} = \hat{Q}^\ast A = \hat{R}$.
\end{proof}

Put differently: the part of $\spn(\vec{a}_1, \dots, \vec{a}_k)$ orthogonal to $\spn(\vec{a}_1, \dots, \vec{a}_{k-1})$ is one-dimensional, so the direction of $\vec{q}_k$ is completely determined by $A$ up to a unimodular phase, and $r_{kk} > 0$ fixes exactly that phase.

\section{(SUPPLEMENT) Further Topics}

\paragraph{Three QR algorithms.}
\begin{center}
\begin{tabular}{lllll}
  \toprule
  & Factorization & Form of $Q$ & Orthogonality of $Q$ & Dense flops \\
  \midrule
  MGS & reduced & explicit $\hat{Q}$ & loses $\sim\epsm\kappa(A)$ & $2mn^2$ \\
  Householder & full & implicit ($n$ vectors $\vec{v}_k$) & orthogonal by definition & $2mn^2 - \frac{2}{3}n^3$ \\
  Givens & full & implicit (pairs $(c, s)$) & orthogonal by definition & $3mn^2 - n^3$ \\
  \bottomrule
\end{tabular}
\end{center}

\paragraph{Householder as an existence proof.} The Householder algorithm never breaks down: the only case needing care is $\vec{x} = A[k{:}m, k] = \zerovec$, where one simply takes $F_k = I$. It therefore produces a full QR factorization of \emph{any} $A$, rank-deficient or not, which is a shorter existence proof than the one based on Gram--Schmidt. Trefethen and Bau prove existence with Gram--Schmidt in Lecture 7, while the algorithm of their Lecture 10 implicitly gives the second proof~\cite{TB}.

\paragraph{Uniqueness via Cholesky.} If $A = \hat{Q}\hat{R}$, then
\begin{equation}
  A^\ast A = \hat{R}^\ast\hat{Q}^\ast\hat{Q}\hat{R} = \hat{R}^\ast\hat{R}.
\end{equation}
For $A$ of full column rank, $A^\ast A$ is Hermitian positive definite, and its Cholesky factorization with positive diagonal is unique; hence $\hat{R}$ is unique, and so is $\hat{Q} = A\hat{R}^{-1}$. This route is not suitable for \emph{computing} QR, since forming $A^\ast A$ squares the condition number (\Cref{ch:ls}).

\paragraph{Computing $c$ and $s$ without overflow.} Evaluating $\sqrt{A_{ii}^2 + A_{ji}^2}$ directly can overflow for large entries and underflow for small ones. Implementations scale by the larger entry first: if $\abs{A_{ji}} > \abs{A_{ii}}$, set $\tau = A_{ii}/A_{ji}$, $s = 1/\sqrt{1 + \tau^2}$, $c = s\tau$, and symmetrically otherwise. This gives $(c, s)$ up to a common sign, which still zeros the entry. The function \texttt{hypot} available in many languages works this way. Demmel also describes how to store a rotation in a single number so that it fits into the zeroed entry~\cite[\S3.4.2]{Demmel}.

\paragraph{Other advantages of Givens rotations.}
\begin{itemize}
  \item \emph{Parallelism}: rotations acting on disjoint pairs of rows can be applied simultaneously.
  \item \emph{Updating}: when a row is appended to an already factored $A$, the new row can be eliminated with $n$ Givens rotations instead of refactoring, which is common in recursive least squares.
  \item QR of Hessenberg matrices is the core step of the QR eigenvalue algorithm~\cite[Lectures 28--29]{TB}.
\end{itemize}

\paragraph{Correspondence with Trefethen and Bau.} Full and reduced QR and their existence and uniqueness are Theorems 7.1 and 7.2 of \cite{TB}; the two implicit algorithms are Algorithms 10.2 and 10.3; Givens rotations appear in Exercise 10.4.


\begin{thebibliography}{9}
\bibitem{TB} L.~N. Trefethen and D.~Bau III. \emph{Numerical Linear Algebra}. SIAM, 1997.
\bibitem{Demmel} J.~W. Demmel. \emph{Applied Numerical Linear Algebra}. SIAM, 1997.
\bibitem{GVL} G.~H. Golub and C.~F. Van Loan. \emph{Matrix Computations}, 3rd edition. Johns Hopkins University Press, 1996.
\bibitem{Overton} M.~L. Overton. \emph{Numerical Computing with IEEE Floating Point Arithmetic}. SIAM, 2001.
\bibitem{Higham} N.~J. Higham. \emph{Accuracy and Stability of Numerical Algorithms}, 2nd edition. SIAM, 2002.
\end{thebibliography}

\end{document}
